% Leçon 31 — Vocabulaire et compréhension de texte : économie — Chinois 1ère A — Cours signé OnBuch+
% Programme officiel MINESEC. Compiler avec : tectonic ZHA31-vocabulaire-economie.tex
\newcommand{\DOCMATIERE}{Chinois}
\newcommand{\DOCNIVEAU}{1ère A}
\newcommand{\DOCMODULE}{Module 4 : Activités économiques et monde du travail}
\newcommand{\DOCLECON}{ZHA31}
\newcommand{\DOCTITRE}{Leçon 31 — Vocabulaire et compréhension de texte : économie}
% OnBuch+ — préambule « cours signés » ENRICHI (compilé avec tectonic / XeLaTeX)
% Système visuel « Pop » d'OnBuch+ : fond crème (jamais blanc), bordures encre
% épaisses, ombres « dures » décalées, titres Archivo Black en capitales.
% Version MATHÉMATIQUES : graphes (pgfplots/tikz), boîtes pédagogiques
% (définition, propriété, méthode, attention, le savais-tu, exercice résolu,
% exercice, TP, bilan), page de garde et signature OnBuch+ ancrée en bas.
%
% Macros attendues dans main.tex AVANT \input{preamble} :
%   \DOCMATIERE  \DOCNIVEAU  \DOCMODULE  \DOCLECON (numéro)  \DOCTITRE
\documentclass[11pt]{article}

\usepackage{fontspec}
\usepackage[french]{babel} % français = langue principale ; le chinois via \zh{..} / textebox (xeCJK)
\usepackage{xeCJK}
\usepackage{pifont} % \ding{51} \ding{55} utilisés par les modèles
\usepackage[a4paper,margin=2cm,top=3.4cm,bottom=2.8cm,headheight=1.4cm,footskip=1.3cm]{geometry}
\usepackage{amsmath,amssymb,amsfonts}
\usepackage{mathtools} % \xmapsto, \coloneqq... souvent utilisés par les modèles
\usepackage{cancel} % simplifications barrées (bilans de forces)
% \abs{z} (module, valeur absolue) et \norm{u} (norme), utilisés par les modèles
\providecommand{\abs}{}\let\abs\relax\DeclarePairedDelimiter{\abs}{\lvert}{\rvert}
\providecommand{\norm}{}\let\norm\relax\DeclarePairedDelimiter{\norm}{\lVert}{\rVert}
\usepackage[table]{xcolor} % [table] : fournit aussi \rowcolors (alternance de lignes)
\usepackage{pagecolor}
\usepackage{tikz}
\usetikzlibrary{arrows.meta,positioning,calc,shapes.geometric,shapes.misc,shapes.arrows,decorations.pathmorphing,decorations.pathreplacing,decorations.markings,patterns,shadows,mindmap,trees,fit,backgrounds,matrix,intersections,angles,quotes}
\usepackage{pgfplots}
\usepackage[version=4]{mhchem} % équations nucléaires \ce{^{235}_{92}U}
\usepackage[european,siunitx]{circuitikz} % schémas électriques
\pgfplotsset{compat=1.18}
\usepgfplotslibrary{fillbetween,groupplots} % aires entre courbes (calcul intégral)
\usepackage{enumitem}
\usepackage{fancyhdr}
% [most] charge toutes les bibliothèques tcolorbox (listings, minted, raster,
% documentation...) et gonfle inutilement la mémoire TeX sur un document long
% (~300 boîtes) : on ne charge que ce qui est réellement utilisé.
\usepackage[skins,breakable]{tcolorbox}
\usepackage{graphicx}
\usepackage{colortbl}
\usepackage{booktabs}
\usepackage{tabularx}
\usepackage{diagbox} % case d'en-tête diagonale des tableaux à double entrée
% Colonnes extensibles centrée / gauche / droite (C, L, R), souvent utilisées par les modèles
\newcolumntype{C}{>{\centering\arraybackslash}X}
\newcolumntype{L}{>{\raggedright\arraybackslash}X}
\newcolumntype{R}{>{\raggedleft\arraybackslash}X}
\usepackage{array}
\usepackage{multicol}
\usepackage{multirow}
\usepackage{siunitx}
% parse-numbers=false : accepte aussi \SI{2,5 \times 10^{-3}}{...}
\sisetup{locale=FR,output-decimal-marker={,},group-minimum-digits=4,parse-numbers=false}
\DeclareSIUnit\ohm{\ensuremath{\symup{\Omega}}} % U+2126 absent de la police
\DeclareSIUnit\division{div} % réglages d'oscilloscope (ms/div, V/div)
\DeclareSIUnit\tr{tr} % tours (vitesse de rotation, ex. \SI{1500}{\tr\per\minute})
\DeclareSIUnit\molar{M} % concentration molaire (ex. \SI{3}{\molar}, \SI{1}{\micro\molar})
\DeclareSIUnit\an{an} % année (le modèle confond parfois avec \year, un compteur TeX) — ex. \SI{5}{\kilo\meter\per\an}
\DeclareSIUnit\FCFA{FCFA} % franc CFA (ex. \SI{500}{\FCFA\per\kilo\gram})
\DeclareSIUnit\degreeBrix{\textdegree Brix} % teneur en sucre d'un sirop/jus (réfractométrie)
\DeclareSIUnit\month{mois} % le modèle utilise parfois \month, un compteur TeX, comme unité de durée
\usepackage{qrcode}
% \degree : commande fréquemment utilisée par les modèles pour un angle ou une
% température en dehors de \SI{}{} (non fournie par les paquets chargés).
\providecommand{\degree}{^{\circ}}
% \qed (fin de démonstration), utilisé par les modèles sans amsthm
\providecommand{\qed}{\hfill$\square$}
% \barre : double barre des valeurs interdites dans un tableau de variations (tkz-tab)
\providecommand{\barre}{\ensuremath{\|}}
% environnement proof (amsthm non chargé) utilisé par les modèles
\ifdefined\proof\else\newenvironment{proof}[1][Démonstration]{\par\noindent\textit{#1.}\ }{\qed\par}\fi
\usepackage{lastpage}
\usepackage{hyperref}
\hypersetup{hidelinks}

% ============================================================
% Palette « Pop » OnBuch+ — mêmes hex que la classe P (plus_ui.dart)
% ============================================================
\definecolor{poporange}{HTML}{F59321}
\definecolor{popdark}{HTML}{E07A0C}
\definecolor{popcream}{HTML}{FFF6E8}
\definecolor{popink}{HTML}{1C1714}
\definecolor{poppurple}{HTML}{7A5AE0}
\definecolor{popgreen}{HTML}{2E9E6A}
\definecolor{poppink}{HTML}{E0457B}
\definecolor{popblue}{HTML}{2F7DE1}
\definecolor{popgold}{HTML}{C98A12}
\definecolor{popmuted}{HTML}{8A7F76}
\definecolor{popline}{HTML}{EADFD2}
\definecolor{popgray}{HTML}{8A7F76}
\colorlet{popgrey}{popgray}
% Alias tolérants (le modèle nomme parfois les couleurs différemment)
\colorlet{popwhite}{white}
\colorlet{popblack}{popink}
\colorlet{popred}{poppink}
\colorlet{popyellow}{popgold}
% Teintes claires (fonds de tableaux/encadrés) — le modèle nomme parfois
% les couleurs avec un suffixe L (ex. popblueL) sans les avoir définies.
\colorlet{poporangeL}{poporange!20}
\colorlet{popdarkL}{popdark!20}
\colorlet{poppurpleL}{poppurple!20}
\colorlet{popgreenL}{popgreen!20}
\colorlet{poppinkL}{poppink!20}
\colorlet{popblueL}{popblue!20}
\colorlet{popgoldL}{popgold!20}
\colorlet{popredL}{popred!20}
\colorlet{popgrayL}{popgray!20}
% Teintes claires pour fonds de tableaux / zones
\colorlet{popblueL}{popblue!10!white}
\colorlet{poporangeL}{poporange!14!white}
\colorlet{popgreenL}{popgreen!12!white}
\colorlet{poppurpleL}{poppurple!12!white}
\colorlet{poppinkL}{poppink!12!white}
\colorlet{popgoldL}{popgold!14!white}

\pagecolor{popcream}

% ============================================================
% Typographie
% ============================================================
\defaultfontfeatures{Ligatures=TeX}
\setmainfont{PlusJakartaSans-Medium.ttf}[Path=../../../fonts/,BoldFont=PlusJakartaSans-Bold.ttf]
% Symboles, API et emojis absents de la police principale : repli sur DejaVu Sans (caractères actifs)
\newfontface\symfont{DejaVuSans.ttf}[Path=../../../fonts/]
\makeatletter
\newcommand\symactive[1]{\catcode#1=13 \begingroup\lccode`\~=#1 \lowercase{\endgroup\def~}{{\symfont\char#1}}}
\newcommand\symblank[1]{\catcode#1=13 \begingroup\lccode`\~=#1 \lowercase{\endgroup\def~}{}}
\newcommand\symhyph[1]{\catcode#1=13 \begingroup\lccode`\~=#1 \lowercase{\endgroup\def~}{\mbox{-}}}
\newcount\symc
\newcommand\symrange[2]{\symc=#1 \@whilenum\symc<#2 \do{\expandafter\symactive\expandafter{\the\symc}\advance\symc by 1 }}
\newcommand\symblankrange[2]{\symc=#1 \@whilenum\symc<#2 \do{\expandafter\symblank\expandafter{\the\symc}\advance\symc by 1 }}
\makeatother
\symrange{"0250}{"02FF}\symactive{"2713}\symactive{"2714}\symactive{"2717}\symactive{"2718}\symactive{"2610}\symactive{"2611}\symactive{"2612}
\symrange{"2600}{"27BF}
\catcode"2705=13 \begingroup\lccode`\~="2705 \lowercase{\endgroup\def~}{{\symfont\char"2713}}\symblank{"26BD}\symblank{"26A1}
\symrange{"2460}{"257F}\symactive{"223C}\symactive{"2248}\symactive{"2260}
\symactive{"01F8}\symactive{"01F9}\symactive{"1E3E}\symactive{"1E3F}
\symhyph{"2011}\symhyph{"2E1A}\symblankrange{"1F300}{"1FAFF}\symblank{"FE0F}
% Caractères chinois : WenQuanYi Zen Hei (sous-ensemble GB2312 + pinyin), gras/italique simulés
\setCJKmainfont{WenQuanYiZenHei-subset.ttf}[Path=../../../fonts/,AutoFakeBold=3,AutoFakeSlant=0.2]
\xeCJKsetup{CJKspace=false}
% Pinyin : voyelles à caron (ǎ ǐ ǒ ǔ ǖ ǘ ǚ ǜ) absentes de la police principale -> caractères actifs
\newfontface\pyfont{WenQuanYiZenHei-subset.ttf}[Path=../../../fonts/]
\newcommand\pyactive[1]{\catcode#1=13 \begingroup\lccode`\~=#1 \lowercase{\endgroup\def~}{{\pyfont\char#1}}}
\pyactive{"01CD}\pyactive{"01CE}\pyactive{"01CF}\pyactive{"01D0}\pyactive{"01D1}\pyactive{"01D2}\pyactive{"01D3}\pyactive{"01D4}\pyactive{"01D5}\pyactive{"01D6}\pyactive{"01D7}\pyactive{"01D8}\pyactive{"01D9}\pyactive{"01DA}\pyactive{"01DB}\pyactive{"01DC}
% Maths en sans-serif (Fira Math) avec chiffres et lettres droites de Plus
% Jakarta, pour que les formules se fondent dans le texte.
\usepackage[math-style=ISO,bold-style=ISO]{unicode-math}
\setmathfont{FiraMath-Regular.otf}[Path=../../../fonts/]
\setmathfont{PlusJakartaSans-Medium.ttf}[Path=../../../fonts/,range={up/num,up/{Latin,latin},\mathup}]
\usepackage{icomma} % virgule décimale sans espace en mode math
% Caractères mathématiques Unicode absents des polices de texte (Plus Jakarta,
% Archivo Black) : rendus par leur équivalent mathématique, en texte comme en
% formule — sinon ils s'affichent en carré vide (ex. « Arithmétique dans ℤ »).
\usepackage{newunicodechar}
\newunicodechar{ℝ}{\ensuremath{\mathbb{R}}}
\newunicodechar{ℤ}{\ensuremath{\mathbb{Z}}}
\newunicodechar{ℂ}{\ensuremath{\mathbb{C}}}
\newunicodechar{ℕ}{\ensuremath{\mathbb{N}}}
\newunicodechar{ℚ}{\ensuremath{\mathbb{Q}}}
\newunicodechar{𝔻}{\ensuremath{\mathbb{D}}}
\newunicodechar{α}{\ensuremath{\alpha}}
\newunicodechar{β}{\ensuremath{\beta}}
\newunicodechar{γ}{\ensuremath{\gamma}}
\newunicodechar{δ}{\ensuremath{\delta}}
\newunicodechar{ε}{\ensuremath{\varepsilon}}
\newunicodechar{θ}{\ensuremath{\theta}}
\newunicodechar{λ}{\ensuremath{\lambda}}
\newunicodechar{μ}{\ensuremath{\mu}}
\newunicodechar{σ}{\ensuremath{\sigma}}
\newunicodechar{τ}{\ensuremath{\tau}}
\newunicodechar{φ}{\ensuremath{\varphi}}
\newunicodechar{ψ}{\ensuremath{\psi}}
\newunicodechar{ω}{\ensuremath{\omega}}
\newunicodechar{Σ}{\ensuremath{\Sigma}}
% majuscules produites par \MakeUppercase dans les titres (α-aminés -> Α)
\newunicodechar{Α}{\ensuremath{\alpha}}
\newunicodechar{Β}{\ensuremath{\beta}}
\newunicodechar{Γ}{\ensuremath{\Gamma}}
\newunicodechar{Δ}{\ensuremath{\Delta}}
\newunicodechar{Ε}{\ensuremath{\varepsilon}}
\newunicodechar{Θ}{\ensuremath{\Theta}}
\newunicodechar{Λ}{\ensuremath{\Lambda}}
\newunicodechar{Μ}{\ensuremath{\mu}}
\newunicodechar{Τ}{\ensuremath{\tau}}
\newunicodechar{Φ}{\ensuremath{\Phi}}
\newunicodechar{Ψ}{\ensuremath{\Psi}}
\newunicodechar{Ω}{\ensuremath{\Omega}}
\newunicodechar{↦}{\ensuremath{\mapsto}}
\newunicodechar{∈}{\ensuremath{\in}}
\newunicodechar{∉}{\ensuremath{\notin}}
\newunicodechar{∧}{\ensuremath{\wedge}}
\newunicodechar{⇔}{\ensuremath{\Leftrightarrow}}
\newunicodechar{⟺}{\ensuremath{\Longleftrightarrow}}
\newunicodechar{⇒}{\ensuremath{\Rightarrow}}
\newunicodechar{⟹}{\ensuremath{\Longrightarrow}}
\newunicodechar{∘}{\ensuremath{\circ}}
\newunicodechar{∪}{\ensuremath{\cup}}
\newunicodechar{∩}{\ensuremath{\cap}}
\newunicodechar{≡}{\ensuremath{\equiv}}
\newunicodechar{⊂}{\ensuremath{\subset}}
\newunicodechar{⊥}{\ensuremath{\perp}}
\newunicodechar{∃}{\ensuremath{\exists}}
\newunicodechar{∀}{\ensuremath{\forall}}
\newunicodechar{∛}{\ensuremath{\sqrt[3]{\,}}}
\newunicodechar{∜}{\ensuremath{\sqrt[4]{\,}}}
\newunicodechar{⃗}{\ensuremath{{}^{\to}}}
\newunicodechar{ⁿ}{\ifmmode^{n}\else\textsuperscript{\ensuremath{n}}\fi}
\newunicodechar{ˣ}{\ifmmode^{x}\else\textsuperscript{\ensuremath{x}}\fi}
\newunicodechar{ᵇ}{\ifmmode^{b}\else\textsuperscript{\ensuremath{b}}\fi}
\newunicodechar{ʳ}{\ifmmode^{r}\else\textsuperscript{\ensuremath{r}}\fi}
\newunicodechar{ᵘ}{\ifmmode^{u}\else\textsuperscript{\ensuremath{u}}\fi}
\newunicodechar{ʸ}{\ifmmode^{y}\else\textsuperscript{\ensuremath{y}}\fi}
\newunicodechar{ᵃ}{\ifmmode^{a}\else\textsuperscript{\ensuremath{a}}\fi}
\newunicodechar{ᵉ}{\ifmmode^{e}\else\textsuperscript{\ensuremath{e}}\fi}
\newunicodechar{ᵏ}{\ifmmode^{k}\else\textsuperscript{\ensuremath{k}}\fi}
\newunicodechar{ᵖ}{\ifmmode^{p}\else\textsuperscript{\ensuremath{p}}\fi}
\newunicodechar{ᴺ}{\ifmmode^{N}\else\textsuperscript{\ensuremath{N}}\fi}
\newunicodechar{ᶜ}{\ifmmode^{c}\else\textsuperscript{\ensuremath{c}}\fi}
\newunicodechar{ʰ}{\ifmmode^{h}\else\textsuperscript{\ensuremath{h}}\fi}
\newunicodechar{ᵠ}{\ifmmode^{q}\else\textsuperscript{\ensuremath{q}}\fi}
\newunicodechar{ₙ}{\ifmmode_{n}\else\textsubscript{\ensuremath{n}}\fi}
\newunicodechar{ₐ}{\ifmmode_{a}\else\textsubscript{\ensuremath{a}}\fi}
\newunicodechar{ᵢ}{\ifmmode_{i}\else\textsubscript{\ensuremath{i}}\fi}
\newunicodechar{ᵣ}{\ifmmode_{r}\else\textsubscript{\ensuremath{r}}\fi}
\newunicodechar{ₖ}{\ifmmode_{k}\else\textsubscript{\ensuremath{k}}\fi}
\newunicodechar{ᵦ}{\ifmmode_{\beta}\else\textsubscript{\ensuremath{\beta}}\fi}
\newunicodechar{ₓ}{\ifmmode_{x}\else\textsubscript{\ensuremath{x}}\fi}
\newfontfamily\popdisplay{ArchivoBlack-Regular.ttf}[Path=../../../fonts/]
\newfontfamily\poplogo{SpaceGrotesk-Bold.ttf}[Path=../../../fonts/]
\newfontfamily\popsemi{PlusJakartaSans-SemiBold.ttf}[Path=../../../fonts/]
\newfontfamily\popxbold{PlusJakartaSans-ExtraBold.ttf}[Path=../../../fonts/]
\newfontfamily\popxboldls{PlusJakartaSans-ExtraBold.ttf}[Path=../../../fonts/,LetterSpace=3.0]

\setlist[enumerate]{leftmargin=1.7em,itemsep=5pt,topsep=4pt}
\setlist[itemize]{leftmargin=1.7em,itemsep=4pt,topsep=4pt,label={\color{poporange}\textbullet}}
\setlength{\parindent}{0pt}
\setlength{\parskip}{5pt}
\renewcommand{\arraystretch}{1.25}

% pgfplots : style Pop par défaut
\pgfplotsset{
  every axis/.append style={axis line style={popink,line width=1pt},tick style={popink},
    grid style={popline,line width=0.5pt},label style={font=\small},tick label style={font=\footnotesize}},
  popcurve/.style={poporange,line width=1.8pt,no markers,smooth},
}
% Même style disponible dans un \draw TikZ simple (hors pgfplots)
\tikzset{popcurve/.style={poporange,line width=1.8pt}}
\tikzset{popgrid/.style={popline,very thin}}
\colorlet{popcurve}{poporange} % le nom du style est parfois utilisé comme couleur

% ============================================================
% Pill / squiggle
% ============================================================
\newcommand{\poppill}[3]{%
  \tikz[baseline=-0.55ex]\node[draw=popink,line width=1.2pt,rounded corners=2.6pt,%
    fill=#2,inner xsep=6.5pt,inner ysep=3pt]{%
    \popxboldls\fontsize{7}{7}\selectfont\color{#3}\MakeUppercase{#1}};%
}
\newcommand{\popsquiggle}[1]{%
  \tikz[baseline=-0.4ex]{
    \draw[#1,line width=2.6pt,line cap=round]
      plot [smooth] coordinates {(0,0.09) (0.34,0.01) (0.68,0.21) (1.02,0.01) (1.36,0.10)};
  }%
}
% Mot-clé surligné (marqueur orange sous le texte)
\newcommand{\cle}[1]{{\popxbold\color{popdark}#1}}

% ============================================================
% En-tête / pied
% ============================================================
\pagestyle{fancy}
\fancyhf{}
\renewcommand{\headrulewidth}{0pt}
\renewcommand{\footrulewidth}{0pt}
\lhead{\raisebox{-0.15ex}{\poplogo\fontsize{13}{13}\selectfont\color{popink}OnBuch\color{poporange}+}}
\rhead{\poppill{\DOCMATIERE\ \textbullet\ \DOCNIVEAU\ \textbullet\ Leçon \DOCLECON}{white}{popink}}
\lfoot{\popxbold\fontsize{7.6}{7.6}\selectfont\color{popmuted}\MakeUppercase{Cours signé OnBuch+}\ \textbullet\ {\popsemi\DOCTITRE}}
\rfoot{\tikz[baseline=-0.6ex]\node[rounded corners=8.5pt,fill=popink,minimum height=17pt,minimum width=17pt,inner xsep=5pt,inner sep=0]{\popxbold\fontsize{8}{8}\selectfont\color{popcream}\thepage\,/\,\pageref*{LastPage}};}

% ============================================================
% Titres
% ============================================================
\newcommand{\chaptitre}[2]{%
  \begin{tcolorbox}[enhanced,colback=poporange,colframe=popink,boxrule=2.6pt,arc=11pt,
      drop shadow=popink,left=15pt,right=15pt,top=12pt,bottom=11pt,before skip=3pt,after skip=18pt]
    \poppill{#2}{popink}{popcream}\par\vspace{9pt}
    {\raggedright\hyphenpenalty=10000\exhyphenpenalty=10000\popdisplay\fontsize{22}{24}\selectfont\color{popcream}\MakeUppercase{#1}\par}
  \end{tcolorbox}
}
% Section numérotée : pastille orange avec le numéro + titre Archivo
\newcounter{popsec}
\newcommand{\coursec}[1]{%
  \stepcounter{popsec}\setcounter{popsub}{0}%
  \par\vspace{16pt}\noindent
  \begin{minipage}{\linewidth}
  \tikz[baseline=-0.7ex]\node[draw=popink,line width=1.6pt,fill=poporange,rounded corners=4pt,minimum size=22pt,inner sep=0]{\popdisplay\fontsize{12}{12}\selectfont\color{popcream}\thepopsec};%
  \hspace{8pt}{\popdisplay\fontsize{15}{17}\selectfont\color{popink}\MakeUppercase{#1}}\par
  \vspace{4pt}\noindent{\color{poporange}\rule{36pt}{3.6pt}}
  \end{minipage}\par\nopagebreak\vspace{8pt}
}
\newcounter{popsub}
\newcommand{\courssub}[1]{%
  \stepcounter{popsub}%
  \par\vspace{9pt}\noindent{\popxbold\fontsize{11.5}{13}\selectfont\color{popink}{\color{poporange}\thepopsec.\thepopsub}\hspace{6pt}#1}\par\nopagebreak\vspace{4pt}
}

% ============================================================
% Boîtes pédagogiques — même recette : fond blanc, bordure épaisse colorée,
% ombre dure encre, badge plein. Titre optionnel [#1].
% ============================================================
\tcbset{popbase/.style={breakable,enhanced,colback=white,boxrule=2.3pt,arc=9pt,
  shadow={3pt}{-3pt}{0pt}{popink},left=14pt,right=14pt,top=11pt,bottom=11pt,before skip=11pt,after skip=15pt,
  fontupper=\popsemi\fontsize{10.6}{14.5}\selectfont\color{popink},
  fontlower=\popsemi\fontsize{10.6}{14.5}\selectfont\color{popink}}}
% Coche dessinée (le glyphe ✓ n'existe pas dans les polices du document)
\newcommand{\popcheck}[1]{\tikz[baseline=-0.5ex]\draw[#1,line width=1.6pt,line cap=round,line join=round](-0.13,0.0)--(-0.04,-0.09)--(0.13,0.1);}
\providecommand{\coche}{\popcheck{popgreen}}
\providecommand{\resultat}[1]{\par\smallskip\noindent\fcolorbox{popgreen}{popgreenL}{\parbox{\dimexpr\linewidth-2\fboxsep-2\fboxrule\relax}{\textbf{Résultat.} #1}}\par\smallskip}
\newcommand{\popboxhead}[3]{\poppill{#1}{#2}{white}\ifx\relax#3\relax\else\hspace{6pt}{\popxbold\fontsize{10.6}{12}\selectfont\color{popink}#3}\fi\par\vspace{7pt}}

\newtcolorbox{aretenir}[1][]{popbase,colframe=popgreen,before upper={\popboxhead{À retenir}{popgreen}{#1}}}
% Texte en chinois (document de lecture, dialogue, extrait) : cadre
% d'étude. Un titre optionnel (source, auteur) s'affiche dans l'en-tête.
\newtcolorbox{textebox}[1][]{popbase,colframe=popblue,before upper={\popboxhead{文本 · Texte}{popblue}{#1}},after upper={}}
% Marques ✓ / ✗ (forme correcte / fautive) souvent utilisées par les modèles
\providecommand{\cross}{\textcolor{poppink}{\ensuremath{\times}}}
\providecommand{\xmark}{\cross}\providecommand{\crossmark}{\cross}\providecommand{\cmark}{\ensuremath{\checkmark}}
% Ligne de réponse / justification à compléter (inventée par les modèles)
\providecommand{\justification}[1]{\par\noindent\textit{Justification~:} \dotfill\par}
% \textipa (package tipa non chargé) : la police ne couvre pas l'API -> simple passage
\providecommand{\textipa}[1]{#1}
% Soulignement (ulem non chargé) : \uline, \ul
\providecommand{\uline}[1]{\underline{#1}}\providecommand{\ul}[1]{\underline{#1}}
% \glossaire{\item ...} inventée par les modèles : liste de glossaire
\providecommand{\glossaire}[1]{\begin{itemize}#1\end{itemize}}
% \alert (beamer) inventée par les modèles : texte mis en évidence
\providecommand{\alert}[1]{\textbf{\textcolor{poppink}{#1}}}
% variantes ASCII d'ordinaux écrites en macro
\providecommand{\iere}{\textsuperscript{re}}\providecommand{\ieme}{\textsuperscript{e}}\providecommand{\ere}{\textsuperscript{re}}\providecommand{\eme}{\textsuperscript{e}}\providecommand{\er}{\textsuperscript{er}}\providecommand{\ie}{\textsuperscript{e}}
% Ordinaux français écrits en macro par les modèles : 1\ière, 3\ième
\newcommand{\ière}{\textsuperscript{re}}\newcommand{\ième}{\textsuperscript{e}}
% \exemple (inventée par les modèles) : étiquette « Exemple : »
\providecommand{\exemple}{\textbf{Exemple~:}\ }
% \square utilisable aussi hors mode math (cases à cocher)
\AtBeginDocument{\let\squareMath\square\renewcommand{\square}{\ensuremath{\squareMath}}}
% Mot ou phrase en chinois à l'intérieur d'un paragraphe français
\newcommand{\zh}[1]{\ifmmode\text{#1}\else{#1}\fi}
\let\eng\zh \let\esp\zh \let\all\zh % alias tolérés
% flèches d'intonation inventées par les modèles
\providecommand{\textsearrow}{}\renewcommand{\textsearrow}{\ensuremath{\searrow}}\providecommand{\textnearrow}{}\renewcommand{\textnearrow}{\ensuremath{\nearrow}}
% \fr{..} : mot français (inventé par les modèles)
\providecommand{\fr}[1]{\textit{#1}}
% zhbox : encadré de texte chinois inventé par les modèles
\newenvironment{zhbox}{\begin{quote}}{\end{quote}}
% \courssubsub : titre de niveau 3 inventé par les modèles
\providecommand{\courssubsub}[1]{\par\medskip\noindent\textbf{#1}\par\smallskip}
% \zhbf : texte chinois en gras inventé par les modèles
\providecommand{\zhbf}[1]{\textbf{#1}}
% \vs : « versus » inventé par les modèles
\providecommand{\vs}{\textit{vs}\ }
% \blank : case à compléter inventée par les modèles
\providecommand{\blank}[1][]{\underline{\hspace{1.6em}}}
\newtcolorbox{exemplebox}[1][]{popbase,colframe=poporange,before upper={\popboxhead{Exemple}{poporange}{#1}}}
\newtcolorbox{definition}[1][]{popbase,colframe=popblue,before upper={\popboxhead{Définition}{popblue}{#1}}}
\newtcolorbox{propriete}[1][]{popbase,colframe=poppurple,before upper={\popboxhead{Propriété}{poppurple}{#1}}}
\newtcolorbox{methode}[1][]{popbase,colframe=popgold,before upper={\popboxhead{Méthode}{popgold}{#1}}}
\newtcolorbox{attention}[1][]{popbase,colframe=poppink,before upper={\popboxhead{Attention}{poppink}{#1}}}
\newtcolorbox{experience}[1][]{popbase,colframe=popblue,colback=popblueL,before upper={\popboxhead{Expérience}{popblue}{#1}}}
% « Le savais-tu ? » : carte violette pleine (contexte, culture, Cameroun)
\newtcolorbox{savaistu}[1][]{popbase,colback=poppurple,colframe=popink,
  fontupper=\popsemi\fontsize{10.4}{14.2}\selectfont\color{white},
  before upper={\poppill{Le savais-tu\,?}{white}{poppurple}\ifx\relax#1\relax\else\hspace{6pt}{\popxbold\color{white}#1}\fi\par\vspace{7pt}}}
% Exercice résolu : énoncé en haut, \tcblower puis la solution sur fond crème
\newcounter{popexo}\newcounter{popexr}
\newtcolorbox{exoresolu}[1][]{popbase,colframe=poporange,colbacklower=popcream,
  segmentation style={popink,line width=1.2pt,dashed},
  before upper={\stepcounter{popexr}\popboxhead{Exercice résolu \thepopexr}{poporange}{#1}},
  before lower={\poppill{Solution}{popink}{popcream}\par\vspace{6pt}}}
% Exercice (non résolu en place) — la correction est dans la partie Corrigés
\newtcolorbox{exercice}[1][]{popbase,colframe=popink,
  before upper={\stepcounter{popexo}\popboxhead{Exercice \thepopexo}{popink}{#1}}}
% Corrigé
\newtcolorbox{corrige}[1][]{popbase,colframe=popgreen,colback=popgreenL,
  before upper={\popboxhead{Corrigé}{popgreen}{#1}}}
% Objectifs / prérequis / bilan
\newtcolorbox{objectifs}{popbase,colframe=popink,colback=poporangeL,
  before upper={\popboxhead{Objectifs de la leçon}{popink}{}}}
\newtcolorbox{prerequis}{popbase,colframe=popink,colback=popblueL,
  before upper={\popboxhead{Prérequis}{popblue}{}}}
\newtcolorbox{bilan}{popbase,colframe=popink,colback=white,boxrule=2.8pt,
  before upper={\popboxhead{Fiche bilan}{poporange}{L'essentiel à connaître pour l'examen}}}
% Figure encadrée avec légende
\newtcolorbox{popfigure}[1][]{enhanced,breakable=false,colback=white,colframe=popline,boxrule=1.4pt,arc=8pt,
  left=10pt,right=10pt,top=10pt,bottom=8pt,before skip=10pt,after skip=12pt,halign=center,
  lower separated=false,halign lower=center,
  fontlower=\popsemi\fontsize{9}{11.5}\selectfont\color{popmuted},
  #1}
\newcommand{\legende}[1]{\par\vspace{6pt}{\popsemi\fontsize{9}{11.5}\selectfont\color{popmuted}#1}}

% ============================================================
% Signature OnBuch+
% ============================================================
\newcommand{\poplogoinline}[1]{{\poplogo\fontsize{#1}{#1}\selectfont\color{popink}OnBuch\color{poporange}+}}
\newcommand{\signatureonbuch}{%
  \begin{tcolorbox}[enhanced,colback=popink,colframe=popink,boxrule=2.2pt,arc=10pt,
      drop shadow=poporange,left=14pt,right=14pt,top=10pt,bottom=10pt,before skip=0pt,after skip=0pt]
    \tikz[baseline=-0.7ex]\node[circle,fill=popgreen,draw=popcream,line width=1.3pt,minimum size=22pt,inner sep=0]{\popcheck{white}};%
    \hspace{9pt}\begin{minipage}[c]{0.62\linewidth}
      {\popxbold\fontsize{12}{13}\selectfont\color{popcream}Signé\ \textbullet\ {\poplogo OnBuch\color{poporange}+}}\par\vspace{2pt}
      {\raggedright\popsemi\fontsize{8}{10.5}\selectfont\color{popline}\MakeUppercase{Équipe pédagogique OnBuch+}\\\MakeUppercase{Conforme au programme officiel MINESEC}\par}
    \end{minipage}\hfill
    {\poplogo\fontsize{22}{22}\selectfont\color{popcream}OnBuch\color{poporange}+}
  \end{tcolorbox}%
}

% ============================================================
% Page de garde
% #1 titre · #2 matière · #3 niveau · #4 module · #5 n° leçon · #6 durée ·
% #7 description / objectifs (texte ou itemize)
% ============================================================
\newcommand{\pagegarde}[7]{%
  \thispagestyle{empty}
  \newgeometry{a4paper,margin=1.7cm,top=1.6cm,bottom=1.4cm}%
  \begin{tikzpicture}[remember picture,overlay]
    \fill[poporange!18] ([xshift=-3.2cm,yshift=-3.6cm]current page.north east) circle (4.6cm);
    \fill[poppurple!14] ([xshift=2.4cm,yshift=7.6cm]current page.south west) circle (3.8cm);
  \end{tikzpicture}%
  {\poplogo\fontsize{30}{30}\selectfont\color{popink}OnBuch\color{poporange}+}\hfill
  \raisebox{0.6ex}{\poppill{Cours signé \textbullet\ Premium}{popink}{popcream}}\par
  \vspace{1pt}\popsquiggle{poppurple}\par
  \vspace{8pt}
  \begin{tcolorbox}[enhanced,colback=poporange,colframe=popink,boxrule=2.8pt,arc=13pt,
      drop shadow=popink,left=18pt,right=18pt,top=12pt,bottom=13pt,after skip=10pt]
    \poppill{#2 \textbullet\ #3}{popink}{popcream}\hspace{5pt}\poppill{#4}{white}{popink}\par\vspace{8pt}
    {\popdisplay\fontsize{14}{14}\selectfont\color{popink}LEÇON #5}\par\vspace{4pt}
    {\raggedright\hyphenpenalty=10000\exhyphenpenalty=10000\popdisplay\fontsize{25}{27}\selectfont\color{popcream}\MakeUppercase{#1}\par}
    \vspace{8pt}
    {\popxbold\fontsize{9.5}{9.5}\selectfont\color{popink}\MakeUppercase{Durée conseillée : #6}}
  \end{tcolorbox}
  \begin{tcolorbox}[enhanced,colback=white,colframe=popink,boxrule=2.2pt,arc=10pt,
      drop shadow=popink,left=16pt,right=16pt,top=10pt,bottom=10pt,after skip=8pt]
    \poppill{Dans cette leçon}{poppurple}{white}\par\vspace{5pt}
    \popsemi\fontsize{9.3}{12.6}\selectfont\color{popink}#7
  \end{tcolorbox}
  \noindent\begin{minipage}[t]{0.487\linewidth}
    \begin{tcolorbox}[enhanced,colback=popgreen,colframe=popink,boxrule=2.2pt,arc=10pt,
        drop shadow=popink,left=11pt,right=11pt,top=10pt,bottom=10pt,height=3.1cm,valign=center]
      \tcbox[colback=white,colframe=popink,boxrule=1.3pt,arc=5pt,boxsep=0pt,nobeforeafter,
        left=3pt,right=3pt,top=3pt,bottom=3pt,valign=center]{\qrcode[height=1.9cm]{https://play.google.com/store/apps/details?id=com.luvvix.onbuch&hl=fr}}%
      \hspace{8pt}\begin{minipage}[c]{0.5\linewidth}
        {\popdisplay\fontsize{11}{12.5}\selectfont\color{white}\MakeUppercase{Télécharge OnBuch}}\par\vspace{4pt}
        {\raggedright\popsemi\fontsize{8.4}{11}\selectfont\color{white}Gratuit \textbullet\ Google Play\\Tuteur IA, annales, concours, résultats\par}
      \end{minipage}
    \end{tcolorbox}
  \end{minipage}\hfill
  \begin{minipage}[t]{0.487\linewidth}
    \begin{tcolorbox}[enhanced,colback=poppurple,colframe=popink,boxrule=2.2pt,arc=10pt,
        drop shadow=popink,left=11pt,right=11pt,top=10pt,bottom=10pt,height=3.1cm,valign=center]
      \tikz[baseline=-0.7ex]\node[circle,fill=white,draw=popink,line width=1.3pt,minimum size=1.6cm,inner sep=0]{\popdisplay\fontsize{24}{24}\selectfont\color{poppurple}+};%
      \hspace{8pt}\begin{minipage}[c]{0.6\linewidth}
        {\popdisplay\fontsize{11}{12.5}\selectfont\color{white}\MakeUppercase{Passe à OnBuch+}}\par\vspace{4pt}
        {\raggedright\popsemi\fontsize{8.4}{11}\selectfont\color{white}Tous les cours signés, Tuteur IA illimité, fiches PDF — onglet OnBuch+ de l'app\par}
      \end{minipage}
    \end{tcolorbox}
  \end{minipage}
  \vspace{8pt}
  \popcontenu
  \vfill
  \signatureonbuch
  \restoregeometry
  \newpage
}

% Rangée « Ce cours contient » (page de garde)
\newcommand{\poptile}[3]{% #1 couleur #2 gros texte #3 légende
  \begin{tcolorbox}[enhanced,colback=#1,colframe=popink,boxrule=2pt,arc=9pt,shadow={2.5pt}{-2.5pt}{0pt}{popink},
      left=6pt,right=6pt,top=9pt,bottom=9pt,halign=center,valign=center,height=1.75cm,width=\linewidth,nobeforeafter]
    {\popdisplay\fontsize{12.5}{13}\selectfont\color{white}\MakeUppercase{#2}}\par\vspace{3pt}
    {\centering\popsemi\fontsize{7.8}{9.5}\selectfont\color{white}#3\par}
  \end{tcolorbox}}
\newcommand{\popcontenu}{%
  \par\noindent{\popxboldls\fontsize{7.5}{7.5}\selectfont\color{popmuted}CE COURS SIGNÉ CONTIENT}\par\vspace{7pt}
  \noindent\begin{minipage}[t]{0.235\linewidth}\poptile{popblue}{Cours}{complet, illustré, pas à pas}\end{minipage}\hfill
  \begin{minipage}[t]{0.235\linewidth}\poptile{poporange}{Méthodes}{et exercices résolus}\end{minipage}\hfill
  \begin{minipage}[t]{0.235\linewidth}\poptile{poppink}{Exercices}{type examen + corrigés}\end{minipage}\hfill
  \begin{minipage}[t]{0.235\linewidth}\poptile{popgreen}{Fiche bilan}{l'essentiel à retenir}\end{minipage}\par}

% Sommaire
\newtcolorbox{sommaire}{enhanced,colback=white,colframe=popink,boxrule=2.2pt,arc=10pt,
  drop shadow=popink,left=18pt,right=18pt,top=13pt,bottom=13pt,before skip=6pt,after skip=20pt,
  before upper={\poppill{Sommaire}{poppurple}{white}\par\vspace{9pt}\popsemi\fontsize{10.6}{14.5}\selectfont\color{popink}}}

% Signature de fin de cours — poussée tout en bas de la dernière page
\newcommand{\findecours}{%
  \par\vfill\nopagebreak
  \begin{center}\popsquiggle{poporange}\end{center}\vspace{4pt}
  \signatureonbuch
}
\AtBeginDocument{\def\llbracket{\mathopen{[\mkern-3.5mu[}}\def\rrbracket{\mathclose{]\mkern-3.5mu]}}} % intervalles d'entiers (glyphe ⟦ absent de la police)
% Glyphes absents de Fira Math : redéfinis à partir de glyphes disponibles
\AtBeginDocument{%
  \def\setminus{\mathbin{\backslash}}\let\smallsetminus\setminus
  \def\vdots{\mathord{\vbox{\baselineskip3.2pt\lineskiplimit0pt\kern4pt\hbox{.}\hbox{.}\hbox{.}}}}%
  \def\ddots{\mathinner{\mkern1mu\raise6.5pt\hbox{.}\mkern2mu\raise3.6pt\hbox{.}\mkern2mu\raise0.7pt\hbox{.}\mkern1mu}}%
  \def\triangle{\mathord{\scalebox{0.9}{$\symup{\Delta}$}}}%
  \def\nabla{\mathord{\raisebox{\depth}{\scalebox{1}[-1]{$\symup{\Delta}$}}}}%
  \def\checkmark{\popcheck{popdark}}%
  \def\star{\mathbin{\ast}}\def\bigstar{\mathord{\ast}}%
  \def\diamond{\mathbin{\scalebox{0.6}{$\Diamond$}}}%
  \def\bigcup{\mathop{\mathchoice{\raisebox{-0.3em}{\scalebox{1.7}{$\cup$}}}{\scalebox{1.25}{$\cup$}}{\cup}{\cup}}\displaylimits}%
  \def\bigcap{\mathop{\mathchoice{\raisebox{-0.3em}{\scalebox{1.7}{$\cap$}}}{\scalebox{1.25}{$\cap$}}{\cap}{\cap}}\displaylimits}%
  \def\lesseqqgtr{\mathrel{\lessgtr}}\def\bigoplus{\mathop{\oplus}\displaylimits}%
}
\providecommand{\ssi}{si et seulement si} % abréviation fréquente
\AtBeginDocument{\providecommand{\wideparen}{\overparen}} % arc de cercle

%% FIGFIX-BEGIN v4 — correctifs globaux des figures (docs/onbuch-generation/tools/figfix)
\usepackage{adjustbox}\usepackage{etoolbox}
% toute figure TikZ/pgfplots plus large que la ligne est réduite à la largeur disponible
\BeforeBeginEnvironment{tikzpicture}{\begin{adjustbox}{max width=\linewidth}}
\AfterEndEnvironment{tikzpicture}{\end{adjustbox}}
% légendes pgfplots lisibles par-dessus les courbes
\pgfplotsset{every axis/.append style={legend style={fill=white,fill opacity=0.92,text opacity=1,draw=black!15,inner sep=3pt}}}
% étiquettes d'arêtes (syntaxe edge["..."]) avec fond blanc
\tikzset{every edge quotes/.append style={fill=white,inner sep=1.2pt}}
%% FIGFIX-END
\begin{document}

\pagegarde{Leçon 31 — Vocabulaire et compréhension de texte : économie}{Chinois}{1ère A}{Module 4 : Activités économiques et monde du travail}{ZHA31}{2 h}{Cette leçon fait découvrir aux élèves de 1ère A le vocabulaire chinois de base du domaine économique (économie, commerce, marché, entreprise, prix, développement) à travers un texte support original. Elle développe la lecture, la prononciation, l'écriture des caractères clés et la compréhension écrite.
\vspace{4pt}
\begin{multicols}{2}\small
\begin{itemize}
\item À la fin de cette leçon, je sais lire et comprendre un court texte chinois sur l'économie.
\item À la fin de cette leçon, je sais reconnaître et prononcer correctement (pinyin et tons) au moins 15 mots du thème économique.
\item À la fin de cette leçon, je sais identifier les radicaux (clés) des caractères 经, 济, 商, 市, 钱 et en respecter l'ordre des traits.
\item À la fin de cette leçon, je sais employer le vocabulaire économique dans des phrases simples (collocations).
\item À la fin de cette leçon, je sais répondre à des questions de compréhension sur un texte économique.
\item À la fin de cette leçon, je sais comparer brièvement des réalités économiques du Cameroun et de la Chine avec des mots chinois.
\end{itemize}\end{multicols}}

\begin{sommaire}
\begin{multicols}{2}
\begin{enumerate}
\item Mise en route et découverte du thème
\item Texte support : lecture et compréhension globale
\item Glossaire du thème : économie
\item Clés (radicaux) et ordre des traits
\item Questions de compréhension et exploitation
\item Production guidée et ancrage Cameroun-Chine
\item Activité d'intégration
\item Méthodes et exercices résolus
\item Exercices
\item Corrigés des exercices
\item Fiche bilan
\end{enumerate}
\end{multicols}
\end{sommaire}

\begin{objectifs}
\begin{itemize}
\item À la fin de cette leçon, je sais lire et comprendre un court texte chinois sur l'économie.
\item À la fin de cette leçon, je sais reconnaître et prononcer correctement (pinyin et tons) au moins 15 mots du thème économique.
\item À la fin de cette leçon, je sais identifier les radicaux (clés) des caractères 经, 济, 商, 市, 钱 et en respecter l'ordre des traits.
\item À la fin de cette leçon, je sais employer le vocabulaire économique dans des phrases simples (collocations).
\item À la fin de cette leçon, je sais répondre à des questions de compréhension sur un texte économique.
\item À la fin de cette leçon, je sais comparer brièvement des réalités économiques du Cameroun et de la Chine avec des mots chinois.
\item À la fin de cette leçon, je sais produire un court paragraphe guidé sur l'économie de mon pays.
\end{itemize}
\end{objectifs}

\begin{prerequis}
\begin{itemize}
\item Maîtrise du pinyin et des quatre tons (Seconde)
\item Vocabulaire des nombres et de l'argent (多少钱, 块) vu en Seconde
\item Structures de base : sujet + verbe + complément ; 是, 有, 很
\item Notion de radicaux (clés) et ordre général des traits
\item Vocabulaire des pays et nationalités (中国, 喀麦隆)
\end{itemize}
\end{prerequis}

\newpage

\coursec{Mise en route et découverte du thème}

\courssub{Situation d'accroche : au marché Mokolo}

Il est huit heures du matin au marché Mokolo de Yaoundé. Les étals débordent de tomates, d'oignons, de plantains. Amina, élève en classe de Première, accompagne sa mère. Devant un étal, la négociation commence : « Maman, c'est combien le tas de tomates ? — Mille francs. — Ah, c'est trop cher ! Hier, c'était sept cents ! » La vendeuse sourit : « Ma fille, les tomates sont rares cette semaine, alors le prix monte. » Amina vient de toucher du doigt, sans le savoir, deux grandes lois de l'économie : \cle{l'offre} et \cle{la demande}. Quand un produit est rare et que tout le monde le veut, le prix monte ; quand il est abondant, le prix baisse.

Le soir, son père écoute les informations à la radio. On parle de la croissance économique du Cameroun, des échanges commerciaux avec la Chine, des entreprises qui créent des emplois, du développement du pays. Des mots reviennent sans cesse : « commerce », « marché », « entreprise », « prix », « argent », « développement ».

\textbf{Problématique :} comment dit-on tout cela en chinois ? Comment un Chinois parle-t-il du marché, de l'argent, des prix, du commerce entre son pays et l'Afrique ? Aujourd'hui, nous allons lire un court texte chinois sur l'économie et apprendre le vocabulaire indispensable pour en parler. À la fin de la leçon, tu seras capable de reconnaître, prononcer et employer les mots essentiels du thème, et de comprendre un texte simple sur l'économie.

\courssub{Brainstorming : quels mots associez-vous à « économie » ?}

Avant d'ouvrir le chinois, activons d'abord ce que tu sais déjà en français. Ferme les yeux et pense au mot « économie ». Quels mots te viennent à l'esprit ?

En classe entière, l'enseignant note au tableau toutes les propositions des élèves. Voici les mots qui reviennent le plus souvent :

\begin{itemize}
\item le \cle{marché} : le lieu où l'on achète et où l'on vend (Mokolo, le marché central de Douala...) ;
\item l'\cle{argent} : ce qui sert à payer (les francs CFA, les billets, le mobile money) ;
\item le \cle{commerce} : l'activité d'achat et de vente, entre personnes, entre villes, entre pays ;
\item les \cle{prix} : la valeur en argent d'un produit, qui monte ou qui baisse ;
\item le \cle{travail} : l'activité qui permet de gagner de l'argent ;
\item l'\cle{entreprise} : une organisation qui produit ou vend des biens et des services ;
\item le \cle{développement} : le progrès économique d'un pays.
\end{itemize}

\begin{attention}[Le mot « économie » a plusieurs sens]
En français, « économie » désigne à la fois la \emph{science} qui étudie la production et les échanges (« il étudie l'économie »), le \emph{système} d'un pays (« l'économie camerounaise ») et le fait d'\emph{épargner} (« faire des économies »). En chinois, le mot \zh{经济} (jīngjì) couvre les deux premiers sens ; pour « faire des économies », on utilise un autre verbe, \zh{省钱} (shěng qián, « économiser de l'argent »). Attention donc à ne pas tout traduire par un seul mot !
\end{attention}

\courssub{Flashcards : les quatre mots-clés du jour}

L'enseignant présente maintenant quatre cartes (flashcards). Sur le recto, le caractère chinois ; sur le verso, le pinyin et la traduction. Répète chaque mot à voix haute trois fois, en faisant bien entendre les tons.

\begin{center}
\begin{tabularx}{\linewidth}{|X|X|X|X|}
\hline
\rowcolor{popblueL} Caractères & Pinyin & Nature & Traduction \\
\hline
经济 & jīngjì & nom / adjectif & économie ; économique \\
\hline
市场 & shìchǎng & nom & marché \\
\hline
钱 & qián & nom & argent \\
\hline
价格 & jiàgé & nom & prix \\
\hline
\end{tabularx}
\end{center}

Observons chaque mot de près :

\begin{definition}[经济 jīngjì — économie]
Mot de deux caractères. \zh{经} (jīng, premier ton) évoque l'idée d'« organiser, gérer » ; \zh{济} (jì, quatrième ton) évoque l'idée d'« aider, secourir ». À l'origine, \zh{经济} signifiait « gérer le monde et aider le peuple » — une belle définition de l'économie ! Aujourd'hui, c'est le mot courant pour « économie » et « économique » : \zh{中国经济} (Zhōngguó jīngjì, l'économie chinoise), \zh{经济发展} (jīngjì fāzhǎn, le développement économique).
\end{definition}

\begin{definition}[市场 shìchǎng — marché]
\zh{市} (shì, quatrième ton) signifie « ville, marché » ; \zh{场} (chǎng, troisième ton) signifie « lieu, espace ». Le marché est donc littéralement le « lieu de la ville » où l'on échange. \zh{市场} désigne aussi bien le marché physique (le marché Mokolo) que le marché abstrait des économistes (le marché africain, le marché mondial).
\end{definition}

\begin{definition}[钱 qián — argent]
Un seul caractère, deuxième ton. Regarde sa clé à gauche : \zh{钅}, la clé du \emph{métal} — car autrefois, la monnaie était en métal (pièces de cuivre, d'argent, d'or). C'est un mot de très haute fréquence : \zh{多少钱？} (Duōshao qián ? — Combien ça coûte ?), la phrase la plus utile au marché !
\end{definition}

\begin{definition}[价格 jiàgé — prix]
\zh{价} (jià, quatrième ton) contient la clé de l'\emph{homme} \zh{亻} à gauche : le prix, c'est la valeur que les hommes donnent aux choses. \zh{格} (gé, deuxième ton) évoque une norme, un standard. \zh{价格} est le mot soutenu pour « prix » ; dans la conversation quotidienne, on dit souvent simplement \zh{价钱} (jiàqian).
\end{definition}

\begin{exemplebox}[Premières phrases avec les mots-clés]
\zh{经济很重要。} \emph{Jīngjì hěn zhòngyào.} « L'économie est très importante. »\\
\zh{这个市场很大。} \emph{Zhège shìchǎng hěn dà.} « Ce marché est très grand. »\\
\zh{我没有钱。} \emph{Wǒ méiyǒu qián.} « Je n'ai pas d'argent. »\\
\zh{西红柿的价格很高。} \emph{Xīhóngshì de jiàgé hěn gāo.} « Le prix des tomates est élevé. »
\end{exemplebox}

Remarque la structure de la dernière phrase : \zh{西红柿的价格}, littéralement « tomates + \zh{的} + prix », c'est-à-dire « le prix \emph{des} tomates ». Le petit mot-outil \zh{的} (de) marque la possession ou la relation, comme « de » en français : A \zh{的} B = « le B de A ».

\begin{propriete}[Rappel de grammaire : le mot-outil 的 (de)]
Structure : \textbf{A + 的 + B} = « le B de A » (possession ou relation).\\
\zh{我的妈妈} \emph{wǒ de māma} « ma mère » ; \zh{中国的经济} \emph{Zhōngguó de jīngjì} « l'économie de la Chine » ; \zh{市场的价格} \emph{shìchǎng de jiàgé} « les prix du marché ».\\
Exception : après un pronom personnel suivi d'un nom de parenté ou d'un mot très proche, on peut omettre \zh{的} : \zh{我妈妈} (wǒ māma, « ma mère ») est correct et très courant.
\end{propriete}

\courssub{Première écoute du texte support}

Maintenant, ferme ton livre. L'enseignant va lire à voix haute, deux fois, un court texte chinois. Tu ne vois pas le texte écrit : tu écoutes seulement. Ta mission : reconnaître les quatre mots-clés du jour (\zh{经济}, \zh{市场}, \zh{钱}, \zh{价格}) et tout autre mot que tu connais déjà (\zh{中国}, \zh{人}, \zh{工作}...). Lève la main ou note un trait sur ton cahier chaque fois que tu entends l'un de ces mots.

\begin{textebox}[Phrase d'ouverture du texte — premier contact]
经济很重要。\\
\emph{Jīngjì hěn zhòngyào.}\\
« L'économie est très importante. »
\end{textebox}

Cette phrase, qui ouvre le texte, donne immédiatement le ton : le texte va expliquer \emph{pourquoi} l'économie est importante dans la vie quotidienne, au Cameroun comme en Chine. Le texte complet sera lu, étudié et traduit dans la section suivante.

\begin{methode}[Avant de lire : deviner le thème grâce aux mots connus]
Face à un texte chinois nouveau, ne panique jamais. Suis ces quatre étapes :
\begin{enumerate}
\item \textbf{Écoute ou survole} le texte une première fois sans chercher à tout comprendre.
\item \textbf{Repère} les mots que tu connais déjà : noms de pays, chiffres, mots appris en classe. Entoure-les mentalement.
\item \textbf{Devine le thème général} : si tu reconnais \zh{市场}, \zh{钱} et \zh{价格}, le texte parle très probablement d'économie ou de commerce.
\item \textbf{Formule une hypothèse} en français : « Je pense que ce texte parle de... ». Tu vérifieras ton hypothèse lors de la lecture détaillée.
\end{enumerate}
Cette stratégie, appelée \emph{lecture globale}, est exactement celle que tu utilises en français quand tu regardes le titre et les images d'un article avant de le lire.
\end{methode}

\courssub{Hypothèses des élèves : que contient le texte ?}

Après la double écoute, l'enseignant interroge la classe : « D'après les mots que vous avez reconnus, de quoi parle ce texte ? » Chaque élève formule une hypothèse en français. Voici des hypothèses typiques, toutes acceptables à ce stade :

\begin{itemize}
\item « J'ai entendu \zh{中国} (Zhōngguó, Chine) et \zh{经济} : le texte parle peut-être de l'économie de la Chine. »
\item « J'ai entendu \zh{市场} et \zh{价格} : peut-être qu'on parle des prix au marché. »
\item « J'ai entendu \zh{工作} (gōngzuò, travailler) et \zh{钱} : peut-être que le texte parle du travail et de l'argent. »
\item « J'ai entendu \zh{人} (rén, personnes) plusieurs fois : le texte parle sûrement des gens et de leur vie. »
\end{itemize}

L'enseignant note les hypothèses au tableau, sans les valider encore. Elles seront vérifiées à la lecture détaillée du texte dans la section 2. Cette étape est essentielle : elle transforme l'écoute passive en écoute \emph{active}, avec un but précis.

\begin{exoresolu}[Associer le mot chinois à sa traduction]
Énoncé : voici quatre mots entendus pendant l'écoute — \zh{经济}, \zh{市场}, \zh{钱}, \zh{价格} — et quatre traductions en désordre : argent, prix, économie, marché. Associe chaque mot chinois à sa traduction, puis complète la phrase : \zh{市场的\_\_\_\_很高。} (Le \_\_\_\_ du marché est très élevé.)
\tcblower
Solution détaillée : \zh{经济} (jīngjì) = économie ; \zh{市场} (shìchǎng) = marché ; \zh{钱} (qián) = argent ; \zh{价格} (jiàgé) = prix. Pour la phrase à trous, le sens exige « prix » : on dit « le \emph{prix} est élevé », jamais « l'argent est élevé » ni « l'économie est élevée ». La phrase complète est donc : \zh{市场的价格很高。} \emph{Shìchǎng de jiàgé hěn gāo.} « Les prix du marché sont très élevés. » Astuce de vérification : l'adjectif \zh{高} (gāo, haut/élevé) se combine naturellement avec \zh{价格} ; en chinois comme en français, un prix « monte » (\zh{价格高}) ou « baisse » (\zh{价格低}).
\end{exoresolu}

\begin{attention}[Piège des francophones : 价格 ou 钱 ?]
Les élèves francophones confondent souvent \zh{价格} (jiàgé, le prix) et \zh{钱} (qián, l'argent). Retiens bien : le \zh{价格} est la \emph{valeur affichée} d'un produit (le prix des tomates), tandis que \zh{钱} est le \emph{moyen de payer} (les billets dans ta poche). On dit \zh{价格很高} (le prix est élevé) mais \zh{钱很多} (l'argent est abondant, « il y a beaucoup d'argent »). Ne dis jamais \zh{钱很高} pour « c'est cher » ! Pour « c'est cher », dis \zh{太贵了} (tài guì le) ou \zh{价格很高}.
\end{attention}

\courssub{Carte mentale : le champ lexical de l'économie}

Pour mémoriser durablement, organisons les mots en réseau. Au centre, le mot vedette \zh{经济} ; autour, les mots de la même famille que nous rencontrerons dans le texte : \zh{市场} (marché), \zh{钱} (argent), \zh{公司} (gōngsī, entreprise), \zh{价格} (prix), \zh{发展} (fāzhǎn, développement), \zh{贸易} (màoyì, commerce/échanges commerciaux).

\begin{popfigure}
\begin{tikzpicture}[mindmap, grow cyclic, every node/.style={concept}, concept color=popblueL,
level 1/.append style={level distance=3.2cm, sibling angle=60},
root concept/.append style={concept color=popblue, font=\Large}]
\node[root concept, text=white, align=center]{经济\\ jīngjì}
child[concept color=poporangeL] { node[align=center]{市场\\ shìchǎng\\ marché} }
child[concept color=popgreenL] { node[align=center]{钱\\ qián\\ argent} }
child[concept color=poppurpleL] { node[align=center]{公司\\ gōngsī\\ entreprise} }
child[concept color=poppinkL] { node[align=center]{价格\\ jiàgé\\ prix} }
child[concept color=popgoldL] { node[align=center]{发展\\ fāzhǎn\\ développement} }
child[concept color=poporangeL] { node[align=center]{贸易\\ màoyì\\ commerce} };
\end{tikzpicture}
\legende{Carte mentale du champ lexical de l'économie : autour de 经济 (jīngjì), six mots associés à mémoriser en réseau.}
\end{popfigure}

Recopie cette carte mentale sur ton cahier : les psychologues de l'apprentissage ont montré qu'on retient mieux un mot quand on le relie à d'autres mots de la même famille, plutôt qu'en liste isolée. Tu pourras ajouter de nouvelles branches au fur et à mesure des leçons du module (\zh{工作} travail, \zh{工人} ouvrier, \zh{银行} banque...).

\courssub{Le savais-tu ? Cameroun et Chine, partenaires économiques}

\begin{savaistu}[La Chine, premier partenaire commercial de l'Afrique]
Le savais-tu ? La Chine est aujourd'hui le premier partenaire commercial du continent africain. Une grande partie des produits vendus au marché Mokolo — chaussures, tissus, appareils électriques, téléphones, casseroles — vient de Chine, importée par des commerçants camerounais qui voyagent parfois jusqu'à Guangzhou ou Yiwu pour acheter leurs marchandises. En sens inverse, le Cameroun exporte vers la Chine du bois, du coton et du cacao. Des entreprises chinoises participent aussi à la construction de routes, de stades et de barrages au Cameroun. Apprendre le chinois, c'est donc aussi se donner un atout concret pour le commerce et le monde du travail de demain : les mots \zh{贸易} (màoyì, commerce) et \zh{经济} (jīngjì, économie) que tu apprends aujourd'hui sont ceux des informations économiques quotidiennes !
\end{savaistu}

\courssub{Objectifs de la leçon}

Pour conclure cette mise en route, voici ce que tu devras savoir faire à la fin de la leçon :

\begin{aretenir}[Objectifs de la leçon 31]
\begin{enumerate}
\item \textbf{Comprendre} un court texte chinois sur le thème de l'économie (lecture globale puis détaillée).
\item \textbf{Reconnaître et prononcer} correctement (pinyin et tons) le vocabulaire du thème : \zh{经济}, \zh{市场}, \zh{钱}, \zh{价格}, \zh{公司}, \zh{发展}, \zh{贸易}...
\item \textbf{Employer} ces mots dans des phrases simples et correctes (ex. : \zh{价格很高}, \zh{经济发展很快}).
\item \textbf{Identifier} les clés (radicaux) et l'ordre des traits de quelques caractères du thème.
\item \textbf{Répondre} à des questions de compréhension sur le texte, en français puis en chinois.
\end{enumerate}
\end{aretenir}

\begin{center}
\begin{tabularx}{\linewidth}{|X|X|X|}
\hline
\rowcolor{popblueL} Français & Chinois + pinyin & Exemple-clé \\
\hline
économie & 经济 jīngjì & 经济很重要。L'économie est importante. \\
\hline
marché & 市场 shìchǎng & 这个市场很大。Ce marché est grand. \\
\hline
argent & 钱 qián & 多少钱？Combien ça coûte ? \\
\hline
prix & 价格 jiàgé & 价格很高。Le prix est élevé. \\
\hline
entreprise & 公司 gōngsī & 这家公司很大。Cette entreprise est grande. \\
\hline
développement & 发展 fāzhǎn & 经济发展很快。Le développement est rapide. \\
\hline
commerce & 贸易 màoyì & 中国和非洲的贸易。Le commerce Chine-Afrique. \\
\hline
\end{tabularx}
\end{center}

\begin{experience}[Jeu d'échauffement : « Au marché ! »]
Par groupes de quatre, jouez une mini-scène de deux minutes : un vendeur, deux clients et un observateur. Le vendeur annonce un prix (\zh{五百块！} Wǔbǎi kuài ! « Cinq cents francs ! »), les clients réagissent (\zh{太贵了！} Tài guì le ! « Trop cher ! » / \zh{便宜一点儿！} Piányi yìdiǎnr ! « Moins cher, un peu ! »), et l'observateur note les mots du thème économie utilisés (\zh{钱}, \zh{价格}, \zh{市场}...). Le groupe qui emploie le plus de mots du jour gagne. Cette activité prépare l'oreille et la bouche avant la lecture du texte.
\end{experience}

Les esprits sont échauffés, les quatre mots-clés sont en place, et tu as déjà formulé tes hypothèses. Passons maintenant à la lecture complète du texte support : à toi de vérifier si tes prédictions étaient justes !

\coursec{Texte support : lecture et compréhension globale}

Après avoir découvert le thème de l'économie et activé nos connaissances sur les échanges entre le Cameroun et la Chine, nous allons maintenant lire un court texte authentique. Ce texte est le cœur de la leçon : c'est en lui que nous allons repérer les mots nouveaux, observer les structures et préparer les questions de compréhension.

\courssub{Le texte support : 中国的经济发展很快}

\begin{textebox}[中国的经济发展很快 — L'économie chinoise se développe vite]
中国的经济发展很快。很多人在市场买东西，也在网上买东西。公司越来越多，工作也越来越多。中国和非洲国家做生意，喀麦隆和中国是好朋友，也是好伙伴。经济的发展让人民的生活越来越好。\\
Zhōngguó de jīngjì fāzhǎn hěn kuài. Hěn duō rén zài shìchǎng mǎi dōngxi, yě zài wǎngshàng mǎi dōngxi. Gōngsī yuèláiyuè duō, gōngzuò yě yuèláiyuè duō. Zhōngguó hé Fēizhōu guójiā zuò shēngyi, Kāmàilóng hé Zhōngguó shì hǎo péngyou, yě shì hǎo huǒbàn. Jīngjì de fāzhǎn ràng rénmín de shēnghuó yuèláiyuè hǎo.\\
\emph{L'économie chinoise se développe très vite. Beaucoup de gens achètent des choses au marché, et aussi sur Internet. Il y a de plus en plus d'entreprises et de plus en plus de travail. La Chine fait du commerce avec les pays africains ; le Cameroun et la Chine sont de bons amis et de bons partenaires. Le développement économique rend la vie du peuple de meilleure en meilleure.}
\end{textebox}

\courssub{Comment lire ce texte : la méthode des trois lectures}

\begin{methode}[Lire un texte chinois en trois fois]
\begin{enumerate}
\item \textbf{Première lecture : le sens global.} Lisez silencieusement sans vous arrêter sur les mots inconnus. Posez-vous une seule question : de quoi parle le texte ?
\item \textbf{Deuxième lecture : les mots nouveaux.} Surlignez les mots du glossaire (\zh{经济}, \zh{发展}, \zh{市场}, \zh{公司}, \zh{做生意}, \zh{伙伴}) et devinez leur sens grâce au contexte avant de regarder la traduction.
\item \textbf{Troisième lecture : les détails.} Relisez phrase par phrase en vue des questions de compréhension : qui ? quoi ? où ? quels pays sont cités ?
\end{enumerate}
\end{methode}

\courssub{Lecture à voix haute : travailler les tons}

La lecture orale se fait en deux étapes :

\begin{itemize}
\item \textbf{Lecture en chœur} : toute la classe lit ensemble, phrase par phrase, en suivant le professeur. Attention particulière aux tons : \zh{jīngjì} (1\textsuperscript{er} ton + 4\textsuperscript{e} ton), \zh{fāzhǎn} (1\textsuperscript{er} ton + 3\textsuperscript{e} ton), \zh{yuèláiyuè} (4\textsuperscript{e} + 2\textsuperscript{e} + 4\textsuperscript{e}).
\item \textbf{Lecture individuelle} : chaque élève lit une phrase à voix haute. Le professeur corrige d'abord les tons, puis la fluidité. Un ton faux peut changer le sens : \zh{mǎi} (acheter, 3\textsuperscript{e} ton) et \zh{mài} (vendre, 4\textsuperscript{e} ton) ne se distinguent que par le ton !
\end{itemize}

\begin{attention}[Le piège des tons : mǎi et mài]
\zh{买} \emph{mǎi} (3\textsuperscript{e} ton) = acheter ; \zh{卖} \emph{mài} (4\textsuperscript{e} ton) = vendre. Dans le texte, on lit \zh{买东西} (mǎi dōngxi, « acheter des choses »). Si vous dites \emph{mài dōngxi}, vous dites « vendre des choses » : le sens est inversé ! Les francophones ont tendance à négliger les tons ; en chinois, le ton fait partie du mot, comme une lettre en français. Observez aussi les caractères : \zh{买} et \zh{卖} se ressemblent, mais \zh{卖} porte en haut l'élément \zh{十} — celui qui « a quelque chose en plus » vend.
\end{attention}

\courssub{Observation : la structure 越来越 (yuèláiyuè)}

Lisons attentivement ces trois phrases du texte :

\begin{exemplebox}[Observer avant d'expliquer]
公司越来越多。 \emph{Gōngsī yuèláiyuè duō.} — Il y a de plus en plus d'entreprises.\\
工作也越来越多。 \emph{Gōngzuò yě yuèláiyuè duō.} — Il y a aussi de plus en plus de travail.\\
人民的生活越来越好。 \emph{Rénmín de shēnghuó yuèláiyuè hǎo.} — La vie du peuple devient de meilleure en meilleure.
\end{exemplebox}

Qu'observez-vous ? La séquence \zh{越来越} revient devant chaque adjectif (\zh{多} « nombreux », \zh{好} « bon »). Elle exprime une progression, exactement comme le français « de plus en plus ».

\begin{propriete}[La structure 越来越 + adjectif]
\textbf{Structure :} Sujet + 越来越 (yuèláiyuè) + adjectif.\\
\textbf{Emploi :} exprime qu'une qualité ou une quantité augmente (ou diminue) progressivement : « de plus en plus... ».\\
\textbf{Règles importantes :}
\begin{itemize}
\item L'adjectif qui suit \zh{越来越} ne prend \textbf{jamais} \zh{很} : on dit \zh{越来越好}, jamais \zh{越来越很好}.
\item On ne met pas \zh{了} à la fin de cette structure pour exprimer la progression en cours.
\item On peut ajouter \zh{也} (yě, « aussi ») avant \zh{越来越} pour coordonner : \zh{工作也越来越多}.
\end{itemize}
\textbf{Comparaison avec le français :} le français répète « plus... plus » ou utilise « de plus en plus » ; le chinois encadre l'adjectif avec 越...越, structure fixe et invariable.
\end{propriete}

\begin{attention}[Erreur fréquente des francophones]
\zh{越来越} (yuèláiyuè) est un bloc figé : ne dites jamais « 很来很去 » ni « 越来很 » ! Autre erreur : ajouter \zh{很} par habitude (\zh{越来越很好} est incorrect). Retenez : 越来越 + adjectif \textbf{seul}.
\end{attention}

\begin{center}
\begin{tabularx}{\linewidth}{|X|X|X|}\hline
\rowcolor{popblueL} Structure & Exemple (caractères + pinyin) & Traduction \\ \hline
越来越 + adj. & 生活越来越好。Shēnghuó yuèláiyuè hǎo. & La vie devient de meilleure en meilleure. \\ \hline
越来越 + adj. & 工作越来越多。Gōngzuò yuèláiyuè duō. & Il y a de plus en plus de travail. \\ \hline
也 + 越来越 + adj. & 汉语也越来越难。Hànyǔ yě yuèláiyuè nán. & Le chinois devient aussi de plus en plus difficile. \\ \hline
越来越 + adj. & 雅温得的汽车越来越多。Yǎwēndé de qìchē yuèláiyuè duō. & Il y a de plus en plus de voitures à Yaoundé. \\ \hline
\end{tabularx}
\end{center}

\begin{popfigure}
\begin{tikzpicture}[node distance=1.6cm, >=Stealth]
\node[draw=popblue, fill=popblueL, rounded corners, thick, align=center, minimum width=2.6cm, minimum height=1.1cm] (eco) {经济\\ \emph{jīngjì}\\ économie};
\node[draw=poporange, fill=poporangeL, rounded corners, thick, align=center, minimum width=2.6cm, minimum height=1.1cm, right=3.2cm of eco] (vie) {生活越来越好\\ \emph{shēnghuó yuèláiyuè hǎo}\\ vie de meilleure en meilleure};
\draw[->, very thick, popgreen] (eco) -- node[above, align=center]{发展 \emph{fāzhǎn}\\ développement} (vie);
\end{tikzpicture}
\legende{Le développement (发展) de l'économie (经济) fait progresser la vie (生活越来越好).}
\end{popfigure}

\courssub{Compréhension globale du texte}

Après la première lecture, répondez mentalement à ces questions :

\begin{itemize}
\item \textbf{De quoi parle le texte ?} Du développement rapide de l'économie chinoise et de ses effets : achats au marché et sur Internet, multiplication des entreprises et des emplois, commerce avec l'Afrique, amélioration de la vie du peuple.
\item \textbf{Quels pays ou régions sont cités ?} La Chine (\zh{中国}), les pays africains (\zh{非洲国家}) et le Cameroun (\zh{喀麦隆}).
\item \textbf{Quelle relation est décrite entre le Cameroun et la Chine ?} Ce sont de bons amis (\zh{好朋友}) et de bons partenaires (\zh{好伙伴}).
\end{itemize}

\begin{exoresolu}[Repérage des mots du glossaire dans le texte]
Énoncé : retrouvez dans le texte les cinq mots suivants et recopiez la phrase qui les contient : 经济, 市场, 公司, 做生意, 伙伴.
\tcblower
Solution détaillée :
\begin{itemize}
\item \zh{经济} (jīngjì, économie) : 中国的经济发展很快。
\item \zh{市场} (shìchǎng, marché) : 很多人在市场买东西。
\item \zh{公司} (gōngsī, entreprise) : 公司越来越多。
\item \zh{做生意} (zuò shēngyi, faire du commerce) : 中国和非洲国家做生意。
\item \zh{伙伴} (huǒbàn, partenaire) : 喀麦隆和中国是好朋友，也是好伙伴。
\end{itemize}
Chaque mot a été repéré dans le texte ; deviner le sens par le contexte avant de consulter la traduction est la bonne démarche.
\end{exoresolu}

\begin{exemplebox}[Deux phrases modèles à retenir]
工作越来越多。 \emph{Gōngzuò yuèláiyuè duō.} — Il y a de plus en plus de travail.\\
生活越来越好。 \emph{Shēnghuó yuèláiyuè hǎo.} — La vie devient de meilleure en meilleure.
\end{exemplebox}

\begin{aretenir}[L'essentiel de cette étape]
\begin{itemize}
\item Un texte se lit trois fois : sens global, mots nouveaux, détails.
\item \zh{越来越} + adjectif = « de plus en plus » ; jamais de \zh{很} après.
\item Le texte présente l'économie chinoise, son développement rapide et le partenariat Cameroun-Chine.
\end{itemize}
\end{aretenir}

Nous sommes maintenant prêts à étudier en détail le glossaire du thème « économie » dans la section suivante.

\coursec{Glossaire du thème : économie}

Dans la section précédente, nous avons lu et compris globalement le texte support sur l'économie et le commerce. Vous avez rencontré des mots comme \zh{经济}, \zh{市场}, \zh{发展} ou encore \zh{做生意}. Il est temps maintenant de nous arrêter sur chacun de ces mots : comment les prononcer avec les bons tons, comment les écrire, comment les employer dans une phrase, et quelles erreurs de francophones il faut absolument éviter. Ce glossaire est votre boîte à outils lexicale pour tout le module « Activités économiques et monde du travail ».

\courssub{Tableau de vocabulaire du thème}

Lisez d'abord le tableau à voix haute, colonne par colonne, en suivant le pinyin. Répétez chaque mot trois fois en veillant aux tons.

\begin{center}
\begin{tabularx}{\linewidth}{|X|X|X|X|}
\hline
\rowcolor{popblueL} Caractères & Pinyin & Nature & Traduction \\
\hline
经济 & jīngjì & n. & économie \\
\hline
发展 & fāzhǎn & v./n. & (se) développer, développement \\
\hline
市场 & shìchǎng & n. & marché \\
\hline
价格 & jiàgé & n. & prix \\
\hline
钱 & qián & n. & argent \\
\hline
公司 & gōngsī & n. & entreprise, société \\
\hline
工作 & gōngzuò & n./v. & travail, travailler \\
\hline
生意 & shēngyi & n. & commerce, affaires \\
\hline
贸易 & màoyì & n. & commerce (international) \\
\hline
买 & mǎi & v. & acheter \\
\hline
卖 & mài & v. & vendre \\
\hline
网上 & wǎngshang & n. & Internet, en ligne \\
\hline
伙伴 & huǒbàn & n. & partenaire \\
\hline
生活 & shēnghuó & n. & vie \\
\hline
人民 & rénmín & n. & peuple \\
\hline
\end{tabularx}
\end{center}

\begin{definition}[Quelques précisions d'emploi]
\begin{itemize}
\item \zh{经济} \emph{jīngjì} désigne l'économie d'un pays ou d'une région : \zh{中国的经济} « l'économie de la Chine ». C'est aussi un adjectif au sens de « économique, pas cher » : \zh{这个饭店很经济。} « Ce restaurant est économique. »
\item \zh{发展} \emph{fāzhǎn} est à la fois verbe et nom : \zh{喀麦隆发展很快。} « Le Cameroun se développe vite » (verbe) ; \zh{经济的发展} « le développement de l'économie » (nom).
\item \zh{生意} \emph{shēngyi} : le second caractère se prononce au \cle{ton neutre} ; c'est le mot de la conversation courante pour « le commerce, les affaires ».
\item \zh{贸易} \emph{màoyì} est plus soutenu : il désigne le commerce officiel, souvent international : \zh{中喀贸易} « le commerce sino-camerounais ».
\item \zh{网上} \emph{wǎngshang} : littéralement « sur le filet », c'est-à-dire « en ligne » : \zh{网上买东西} « acheter en ligne ».
\end{itemize}
\end{definition}

\courssub{Prononciation guidée : les tons à surveiller}

Trois pièges de prononciation attendent le francophone dans ce glossaire. Entraînez-vous en chœur, puis individuellement.

\begin{enumerate}
\item \zh{经济} \emph{jīngjì} : tons \cle{1--4}. La première syllabe est haute et plate (comme une note tenue), la seconde descend brusquement. Ne dites pas \emph{jìngjī} : cela inverserait les tons et brouillerait le sens.
\item \zh{市场} \emph{shìchǎng} : tons \cle{4--3}. Le \emph{sh} se prononce avec la langue rétroflexe (recourbée), différent du « ch » français. Le 3\up{e} ton descend puis remonte.
\item \zh{买} \emph{mǎi} (3\up{e} ton) contre \zh{卖} \emph{mài} (4\up{e} ton) : voir la mise en garde ci-dessous.
\end{enumerate}

\begin{attention}[买 mǎi ou 卖 mài ? Le piège classique]
\zh{买} \emph{mǎi} « acheter » (3\up{e} ton) et \zh{卖} \emph{mài} « vendre » (4\up{e} ton) se ressemblent beaucoup : \zh{卖} n'est que \zh{买} coiffé d'un petit \zh{十} en haut. À l'oral, seul le ton les distingue ; à l'écrit, seul ce petit élément supplémentaire. Moyen mnémotechnique : \cle{« mǎi je prends, mài je donne »} — quand j'achète, je prends la marchandise ; quand je vends, je la donne. Autre astuce visuelle : le vendeur \zh{卖} porte quelque chose « sur la tête » (le \zh{十}), comme les vendeuses de nos marchés qui portent leurs plateaux. Comparez : \zh{我买水果。} \emph{Wǒ mǎi shuǐguǒ.} « J'achète des fruits. » / \zh{她卖水果。} \emph{Tā mài shuǐguǒ.} « Elle vend des fruits. »
\end{attention}

\courssub{Les collocations essentielles}

En chinois, un mot seul ne suffit pas : il faut connaître ses compagnons habituels. Voici les quatre collocations à mémoriser comme des blocs.

\begin{center}
\begin{tabularx}{\linewidth}{|X|X|X|}
\hline
\rowcolor{popgreenL} Collocation & Pinyin & Traduction \\
\hline
发展经济 & fāzhǎn jīngjì & développer l'économie \\
\hline
经济发展 & jīngjì fāzhǎn & le développement économique \\
\hline
做生意 & zuò shēngyi & faire du commerce \\
\hline
买东西 & mǎi dōngxi & faire des achats \\
\hline
\end{tabularx}
\end{center}

Observez bien la différence entre \zh{发展经济} (verbe + objet : « développer l'économie ») et \zh{经济发展} (nom + nom : « le développement économique »). L'ordre des mots change la fonction grammaticale, exactement comme en français « développer l'économie » contre « le développement économique ».

\begin{propriete}[做生意 : expression verbale figée]
\zh{做生意} \emph{zuò shēngyi} est une \cle{expression verbale figée} (verbe + objet inséparables) signifiant « faire du commerce, être commerçant ». Dans la conversation courante, on dit \zh{他做生意。} « Il fait du commerce. » On \textbf{ne dit pas} \zh{做贸易} dans la langue parlée : \zh{贸易} appartient au vocabulaire officiel et écrit (\zh{国际贸易} « commerce international »). Comme toute expression verbe + objet, elle peut être coupée par un complément de durée ou de quantité : \zh{他做了十年生意。} « Il a fait du commerce pendant dix ans. »
\end{propriete}

\begin{exemplebox}[Le glossaire en action]
\zh{我在市场买东西。} \emph{Wǒ zài shìchǎng mǎi dōngxi.} « J'achète des choses au marché. »\\
\zh{他们做生意。} \emph{Tāmen zuò shēngyi.} « Ils font du commerce. »\\
\zh{这家公司的价格不贵。} \emph{Zhè jiā gōngsī de jiàgé bù guì.} « Les prix de cette entreprise ne sont pas chers. »\\
\zh{中国和喀麦隆是贸易伙伴。} \emph{Zhōngguó hé Kāmàilóng shì màoyì huǒbàn.} « La Chine et le Cameroun sont des partenaires commerciaux. »\\
\zh{经济发展了，人民的生活就好了。} \emph{Jīngjì fāzhǎn le, rénmín de shēnghuó jiù hǎo le.} « Quand l'économie se développe, la vie du peuple s'améliore. »\\
\zh{现在很多人在网上买东西。} \emph{Xiànzài hěn duō rén zài wǎngshang mǎi dōngxi.} « Aujourd'hui, beaucoup de gens achètent en ligne. »
\end{exemplebox}

Remarquez le classificateur \zh{家} \emph{jiā} dans \zh{这家公司} : c'est le classificateur des entreprises, des magasins et des restaurants. On dit \zh{一家公司} « une entreprise », jamais \zh{一个公司}.

\courssub{Carte mentale du lexique}

\begin{popfigure}
\begin{tikzpicture}[mindmap, grow cyclic, every node/.style={concept, align=center, font=\small}, root concept/.append style={concept color=popblue, font=\normalsize\bfseries}, level 1/.append style={level distance=3.2cm, sibling angle=72}, level 1 concept/.append style={concept color=poporangeL}]
\node[root concept, align=center]{经济\\ jīngjì\\ économie}
  child { node[align=center]{市场 shìchǎng\\ 价格 jiàgé\\ 钱 qián} }
  child { node[align=center]{买 mǎi\\ 卖 mài\\ 买东西} }
  child { node[align=center]{生意 shēngyi\\ 贸易 màoyì\\ 做生意} }
  child { node[align=center]{公司 gōngsī\\ 工作 gōngzuò\\ 伙伴 huǒbàn} }
  child { node[align=center]{发展 fāzhǎn\\ 生活 shēnghuó\\ 人民 rénmín} };
\end{tikzpicture}
\legende{Carte mentale du vocabulaire de l'économie : le mot central 经济 rayonne vers cinq familles de mots.}
\end{popfigure}

\courssub{Texte de réinvestissement}

Lisez ce court texte qui réemploie le glossaire dans un contexte camerounais, puis répondez mentalement aux questions qui suivent.

\begin{textebox}[Le marché de Mokolo]
雅温得的莫科洛市场很大，也很热闹。\\
Yǎwēndé de Mòkēluò shìchǎng hěn dà, yě hěn rènao.\\
Le marché Mokolo de Yaoundé est grand et très animé.\\
在那里，人们买很多东西：水果、蔬菜、鱼和衣服。\\
Zài nàli, rénmen mǎi hěn duō dōngxi : shuǐguǒ, shūcài, yú hé yīfu.\\
Là-bas, les gens achètent beaucoup de choses : fruits, légumes, poisson et vêtements.\\
商人们每天做生意，价格常常不一样。\\
Shāngrénmen měitiān zuò shēngyi, jiàgé chángcháng bù yíyàng.\\
Les commerçants font du commerce chaque jour, et les prix sont souvent différents.\\
现在，一些公司也在网上卖东西。\\
Xiànzài, yìxiē gōngsī yě zài wǎngshang mài dōngxi.\\
Aujourd'hui, certaines entreprises vendent aussi en ligne.\\
中国是喀麦隆重要的贸易伙伴。\\
Zhōngguó shì Kāmàilóng zhòngyào de màoyì huǒbàn.\\
La Chine est un partenaire commercial important du Cameroun.\\
经济发展了，人民的生活就会更好。\\
Jīngjì fāzhǎn le, rénmín de shēnghuó jiù huì gèng hǎo.\\
Quand l'économie se développe, la vie du peuple devient meilleure.
\end{textebox}

Questions de vérification rapide : 1. Où est le marché Mokolo ? 2. Qu'est-ce que les gens achètent au marché ? 3. Qui vend des choses en ligne ? 4. Quel pays est un partenaire commercial du Cameroun ?

\courssub{Exercice résolu et entraînement}

\begin{exoresolu}[Associer le mot chinois à sa traduction]
Énoncé. Reliez chaque mot chinois à sa traduction française : 1. 市场 ; 2. 公司 ; 3. 价格 ; 4. 伙伴 ; 5. 贸易. Traductions : a. prix ; b. partenaire ; c. marché ; d. commerce international ; e. entreprise.
\tcblower
Solution détaillée. 1. \zh{市场} \emph{shìchǎng} → c. marché (pensez au « champ » \zh{场} où se tient le marché). 2. \zh{公司} \emph{gōngsī} → e. entreprise. 3. \zh{价格} \emph{jiàgé} → a. prix. 4. \zh{伙伴} \emph{huǒbàn} → b. partenaire. 5. \zh{贸易} \emph{màoyì} → d. commerce international. Vérifiez toujours la nature du mot : ici, tous sont des noms, donc la traduction doit être un nom français, pas un verbe.
\end{exoresolu}

\begin{exercice}[À vous de jouer]
1. Traduisez en français : \zh{她在公司工作。} / \zh{我们发展经济。} / \zh{网上买东西很方便。}\\
2. Traduisez en chinois (avec pinyin) : « J'achète des légumes au marché. » / « Ils font du commerce à Douala. »\\
3. Complétez avec \zh{买} ou \zh{卖} : \zh{商人……东西，我……东西。}
\end{exercice}

\begin{corrige}[Exercice « À vous de jouer »]
1. « Elle travaille dans une entreprise. » / « Nous développons l'économie. » / « Acheter en ligne est très pratique. »\\
2. \zh{我在市场买蔬菜。} \emph{Wǒ zài shìchǎng mǎi shūcài.} / \zh{他们在杜阿拉做生意。} \emph{Tāmen zài Dù'ālā zuò shēngyi.}\\
3. \zh{商人卖东西，我买东西。} « Le commerçant vend des choses, j'achète des choses. » — souvenez-vous : mǎi je prends, mài je donne.
\end{corrige}

\begin{savaistu}[D'où vient le mot 经济 ?]
\zh{经济} \emph{jīngjì} vient d'une expression ancienne du chinois classique, \zh{经世济民} « gouverner le monde et aider le peuple ». À l'origine, « économie » désignait donc l'art de bien administrer l'État pour le bien du peuple. C'est seulement à l'époque moderne, sous l'influence du japonais, que le mot a pris son sens actuel d'« économie » au sens financier. Une belle leçon : pour la tradition chinoise, l'économie doit d'abord servir le peuple, \zh{人民} !
\end{savaistu}

\begin{aretenir}[L'essentiel du glossaire]
\begin{itemize}
\item 15 mots clés autour de \zh{经济} : marché, prix, argent, entreprise, commerce, partenaire.
\item Tons critiques : \zh{经济} (1--4), \zh{市场} (4--3), \zh{买} mǎi (3) contre \zh{卖} mài (4) — « mǎi je prends, mài je donne ».
\item Collocations figées : \zh{发展经济}, \zh{经济发展}, \zh{做生意}, \zh{买东西}.
\item \zh{生意} = commerce parlé ; \zh{贸易} = commerce officiel, international.
\item Classificateur des entreprises : \zh{家} → \zh{一家公司} « une entreprise ».
\end{itemize}
\end{aretenir}

\coursec{Clés (radicaux) et ordre des traits}

Dans la section précédente, nous avons découvert le glossaire du thème de l'économie : \zh{经济} (jīngjì, économie), \zh{市场} (shìchǎng, marché), \zh{钱} (qián, argent), \zh{商人} (shāngrén, commerçant)… Mais reconnaître un caractère ne suffit pas : à l'examen comme dans la vie réelle, il faut savoir l'\emph{écrire} correctement. Or l'écriture chinoise obéit à des règles précises : chaque caractère possède une \cle{clé} (radical) qui oriente vers son sens, et un \cle{ordre des traits} immuable. C'est ce que nous allons étudier maintenant, caractère par caractère.

\courssub{Rappel : les règles générales de l'ordre des traits}

\begin{propriete}[Les six règles d'or de l'écriture]
\begin{enumerate}
\item \textbf{De haut en bas} : on écrit d'abord les traits du haut, puis ceux du bas. Exemple : \zh{三} (sān, trois) : les trois traits horizontaux se tracent du haut vers le bas.
\item \textbf{De gauche à droite} : pour les caractères composés de deux parties côte à côte, on écrit d'abord la partie gauche, puis la partie droite. Exemple : \zh{你} (nǐ, tu) : la clé 亻 à gauche, puis 尔 à droite.
\item \textbf{L'extérieur avant l'intérieur} : pour les caractères avec enceinte, on trace d'abord le cadre extérieur, puis le contenu, et on ferme en dernier. Exemple : \zh{国} (guó, pays).
\item \textbf{L'horizontal avant le vertical} quand ils se croisent : \zh{十} (shí, dix) : d'abord 一, puis 丨.
\item \textbf{Le trait central avant les côtés} dans les caractères symétriques : \zh{小} (xiǎo, petit).
\item \textbf{La clé d'abord} : dans la grande majorité des caractères composés, la clé se trace en premier.
\end{enumerate}
\end{propriete}

\begin{attention}[Erreur fréquente des francophones]
En français, nous écrivons les lettres dans n'importe quel ordre tant que le résultat est lisible. En chinois, c'est impossible : un mauvais ordre des traits déforme le caractère, le rend illisible à l'écriture rapide, et empêche d'utiliser un dictionnaire classé par clés et par nombre de traits. À l'examen, l'ordre des traits peut être \emph{noté}. Entraîne-toi toujours en traçant à voix haute : « un, deux, trois… ».
\end{attention}

\courssub{Étude caractère par caractère}

\textbf{1. 经 jīng — « traverser, gérer » (dans 经济 jīngjì, économie)}\par
Clé : \cle{纟} (sī, soie), 3 traits. Nombre total de traits : \textbf{8}.
Ordre : on trace d'abord la clé 纟 à gauche (3 traits : 撇折, 撇折, 提), puis la partie droite 巠 (5 traits : 横撇, 点, 横, 竖, 横).
Règle appliquée : gauche avant droite.

\textbf{2. 济 jì — « aider, franchir » (dans 经济)}\par
Clé : \cle{氵} (shuǐ, eau), 3 traits. Nombre total : \textbf{9}.
Ordre : d'abord les trois gouttes d'eau 氵 à gauche (点, 点, 提), puis la partie 齐 à droite, de haut en bas.
Remarque : la clé de l'eau rappelle le sens ancien de 济 : « traverser une rivière ».

\textbf{3. 商 shāng — « commerce » (dans 商人 shāngrén, commerçant)}\par
Clé : \cle{亠} (tóu, couvercle) ; on le retrouve aussi classé sous 口. Nombre total : \textbf{11}.
Ordre : 点, 横, 点, 撇, 竖, 横折钩, 撇, 点, 竖, 横折, 横.
Règle appliquée : haut avant bas, extérieur avant intérieur (le 口 final se trace en dernier, en bas).

\textbf{4. 市 shì — « marché, ville » (dans 市场 shìchǎng, marché)}\par
Clé : \cle{亠}. Nombre total : \textbf{5}.
Ordre : 点, 横, 竖, 横折钩, 竖.
C'est l'un des caractères les plus fréquents du thème économique : 市场 (shìchǎng, marché), 城市 (chéngshì, ville), 超市 (chāoshì, supermarché).

\textbf{5. 钱 qián — « argent »}\par
Clé : \cle{钅} (jīn, métal), 5 traits. Nombre total : \textbf{10}.
Ordre : d'abord la clé 钅 à gauche (撇, 横, 横, 横, 竖提), puis la partie 戋 à droite (横, 横, 斜钩, 撇, 点).

\begin{methode}[La clé 钅 (métal) : un indice de sens précieux]
La clé 钅 (métal) dans \zh{钱} n'est pas là par hasard : dans la Chine ancienne, la monnaie était frappée en métal (cuivre, bronze). Les radicaux donnent très souvent le \emph{champ sémantique} du caractère. Démarche à suivre face à un caractère inconnu :
\begin{enumerate}
\item Repérer la clé (souvent à gauche ou en haut).
\item Deviner le domaine de sens : 钅 → métal, 氵 → eau, 纟 → soie/textile, 口 → bouche/parole.
\item La partie restante donne souvent la \emph{prononciation} (composante phonétique).
\item Vérifier dans le dictionnaire en cherchant par la clé.
\end{enumerate}
Exemple : \zh{银} (yín, argent métal ; 银行 yínháng, banque) contient aussi 钅.
\end{methode}

\begin{popfigure}
\begin{tikzpicture}[node distance=0.9cm, every node/.style={font=\large}]
\node[draw=popblue, fill=popblueL, rounded corners, minimum width=1.1cm, minimum height=1.1cm] (a1) {纟};
\node[right=of a1] (p1) {\textbf{+}};
\node[draw=poporange, fill=poporangeL, rounded corners, minimum width=1.1cm, minimum height=1.1cm, right=of p1] (b1) {巠};
\node[right=of b1] (p2) {\textbf{=}};
\node[draw=popgreen, fill=popgreenL, rounded corners, minimum width=1.1cm, minimum height=1.1cm, right=of p2] (c1) {经};
\node[below=0.35cm of a1, align=center, text=popblue] {\footnotesize clé soie\\ \footnotesize 3 traits};
\node[below=0.35cm of b1, align=center, text=poporange] {\footnotesize phonétique\\ \footnotesize 5 traits};
\node[below=0.35cm of c1, align=center, text=popgreen] {\footnotesize jīng\\ \footnotesize 8 traits};
\node[draw=popblue, fill=popblueL, rounded corners, minimum width=1.1cm, minimum height=1.1cm, below=1.6cm of a1] (a2) {钅};
\node[right=of a2] (p3) {\textbf{+}};
\node[draw=poporange, fill=poporangeL, rounded corners, minimum width=1.1cm, minimum height=1.1cm, right=of p3] (b2) {戋};
\node[right=of b2] (p4) {\textbf{=}};
\node[draw=popgreen, fill=popgreenL, rounded corners, minimum width=1.1cm, minimum height=1.1cm, right=of p4] (c2) {钱};
\node[below=0.35cm of a2, align=center, text=popblue] {\footnotesize clé métal\\ \footnotesize 5 traits};
\node[below=0.35cm of b2, align=center, text=poporange] {\footnotesize phonétique\\ \footnotesize 5 traits};
\node[below=0.35cm of c2, align=center, text=popgreen] {\footnotesize qián\\ \footnotesize 10 traits};
\draw[-{Stealth}, thick, popdark] (a1.south east) -- (c1.north west);
\draw[-{Stealth}, thick, popdark] (a2.south east) -- (c2.north west);
\end{tikzpicture}
\legende{Décomposition des caractères 经 = 纟 + 巠 et 钱 = 钅 + 戋 : la clé (à gauche) donne le sens, la partie droite suggère le son.}
\end{popfigure}

\courssub{Tableau récapitulatif}

\begin{center}
\begin{tabularx}{\linewidth}{|X|X|X|X|}\hline
\rowcolor{popblueL} Caractère & Pinyin / Traduction & Clé & Nombre de traits \\ \hline
经 & jīng — gérer (经济 : économie) & 纟 (soie) & 8 \\ \hline
济 & jì — aider, franchir & 氵 (eau) & 9 \\ \hline
商 & shāng — commerce & 亠 / 口 & 11 \\ \hline
市 & shì — marché, ville & 亠 & 5 \\ \hline
钱 & qián — argent & 钅 (métal) & 10 \\ \hline
\end{tabularx}
\end{center}

\begin{exemplebox}[Le vocabulaire en action]
\zh{市场} \emph{shìchǎng} — « marché ». Exemple : \zh{喀麦隆的市场很热闹。} \emph{Kāmàilóng de shìchǎng hěn rènao.} « Les marchés du Cameroun sont très animés. »\\
\zh{人民币} \emph{rénmínbì} — « monnaie du peuple » (le yuan). Exemple : \zh{人民币是中国的钱。} \emph{Rénmínbì shì Zhōngguó de qián.} « Le renminbi est la monnaie de la Chine. »\\
\zh{经济} \emph{jīngjì} — « économie ». Exemple : \zh{中国经济很大。} \emph{Zhōngguó jīngjì hěn dà.} « L'économie chinoise est très grande. »
\end{exemplebox}

\begin{attention}[Ne pas confondre 市 et 布]
\zh{市} (shì, marché/ville) et \zh{布} (bù, tissu) se ressemblent beaucoup : un seul trait de différence ! 市 commence par un \emph{point} (点) au sommet, tandis que 布 commence par un trait horizontal suivi d'une 撇. Piège de traduction : \zh{超市} (chāoshì) signifie « supermarché », pas « super-ville ». Vérifie toujours le premier trait avant d'écrire.
\end{attention}

\begin{savaistu}[La soie, fil rouge de l'économie chinoise]
La clé 纟 (soie) apparaît dans de nombreux caractères liés au tissu et aux échanges anciens : \zh{纸} (zhǐ, papier), \zh{红} (hóng, rouge — la couleur de la teinture), \zh{线} (xiàn, fil), \zh{给} (gěi, donner). Ce n'est pas un hasard : la soie fut pendant des siècles le premier produit d'exportation de la Chine, au point que la route commerciale vers l'Occident s'appelle la « Route de la soie » (\zh{丝绸之路} Sīchóu zhī Lù). Aujourd'hui encore, le caractère 经 de 经济 (économie) porte cette mémoire du textile !
\end{savaistu}

\begin{exoresolu}[Identifier clé et ordre des traits]
Énoncé : pour le caractère \zh{钱}, indique : a) sa clé et sa signification ; b) son nombre total de traits ; c) la règle générale d'ordre des traits qu'il illustre.
\tcblower
Solution détaillée : a) La clé de 钱 est 钅, la clé du métal (5 traits) : elle rappelle que la monnaie ancienne était en métal. b) 钱 compte 10 traits au total : 5 pour 钅 et 5 pour 戋. c) Il illustre la règle « de gauche à droite » : on trace d'abord la clé 钅 à gauche, puis la composante phonétique 戋 à droite.
\end{exoresolu}

\begin{exercice}[À ton tour : atelier d'écriture]
\begin{enumerate}
\item Trace chacun des cinq caractères du thème (\zh{经}, \zh{济}, \zh{商}, \zh{市}, \zh{钱}) \textbf{5 fois} sur ton cahier, en comptant les traits à voix haute et en respectant l'ordre étudié.
\item Pour chaque caractère, écris à côté : la clé, son sens et le nombre total de traits.
\item Entoure en rouge la clé dans chaque caractère, puis écris une phrase complète avec \zh{市场} et une avec \zh{钱}.
\end{enumerate}
\end{exercice}

\begin{corrige}[Exercice : atelier d'écriture]
1. Vérification par l'enseignant : l'ordre doit être 纟 puis 巠 pour 经 ; 氵 puis 齐 pour 济 ; du haut vers le bas pour 商 ; 点, 横, 竖, 横折钩, 竖 pour 市 ; 钅 puis 戋 pour 钱.\\
2. 经 : clé 纟 (soie), 8 traits. 济 : clé 氵 (eau), 9 traits. 商 : clé 亠/口, 11 traits. 市 : clé 亠, 5 traits. 钱 : clé 钅 (métal), 10 traits.\\
3. Exemples de phrases attendues : \zh{我去市场买东西。} \emph{Wǒ qù shìchǎng mǎi dōngxi.} « Je vais au marché acheter des choses. » \zh{我没有钱。} \emph{Wǒ méiyǒu qián.} « Je n'ai pas d'argent. »
\end{corrige}

\begin{aretenir}[L'essentiel de la section]
\begin{itemize}
\item Chaque caractère possède une \cle{clé} qui oriente vers le sens : 纟 (soie), 氵 (eau), 钅 (métal), 亠 (couvercle).
\item L'ordre des traits suit des règles fixes : haut → bas, gauche → droite, extérieur → intérieur, clé d'abord.
\item 经 (8 traits), 济 (9), 商 (11), 市 (5), 钱 (10) : apprends-les en les traçant, pas seulement en les regardant.
\item Attention aux caractères proches : 市 (marché) ≠ 布 (tissu).
\end{itemize}
\end{aretenir}

\coursec{Questions de compréhension et exploitation}

\courssub{De l'écriture à la compréhension : transition}

Dans la section précédente, nous avons appris à \emph{écrire} les caractères du thème : la clé de \zh{钱} (qián, argent) est 钅, la clé du métal ; \zh{买} (mǎi, acheter) et \zh{卖} (mài, vendre) renvoient à l'échange. Savoir écrire un caractère, c'est bien ; savoir \emph{répondre à une question} avec ce mot, c'est mieux. Nous allons maintenant vérifier notre compréhension du texte sur l'économie, puis apprendre la méthode générale pour répondre à n'importe quelle question en chinois.

\courssub{Méthode : comment répondre à une question en chinois}

\begin{methode}[Répondre en reprenant les mots de la question]
En chinois, la réponse reprend presque toujours les mots de la question, dans le même ordre. Il suffit de remplacer le mot interrogatif par l'information demandée.
\begin{enumerate}
\item \textbf{Identifier le mot interrogatif} : \zh{怎么样？} (comment ?), \zh{哪儿？} (où ?), \zh{谁？} (qui ?), \zh{什么？} (quoi ?), \zh{哪些？} (lesquels ?).
\item \textbf{Recopier la phrase de la question} en supprimant le mot interrogatif.
\item \textbf{Mettre la réponse à la place du mot interrogatif}, sans changer l'ordre des mots.
\end{enumerate}
Exemple : \zh{人们在哪儿买东西？} → on remplace \zh{哪儿} par \zh{市场} → \zh{人们在市场买东西。}
\end{methode}

\begin{attention}[La place de 怎么样 : un piège pour les francophones]
En français, le mot interrogatif se place en tête : « \emph{Comment} se développe l'économie ? ». En chinois, \zh{怎么样} reste à la place de la réponse, c'est-à-dire \textbf{en fin de phrase} : \zh{中国的经济发展怎么样？} (littéralement : « L'économie chinoise se développe comment ? »). Ne dites jamais \zh{怎么样中国的经济发展？} : c'est un calque du français, incorrect en chinois. Même principe pour \zh{哪儿} : \zh{人们在哪儿买东西？} et non \zh{哪儿人们买东西？}. Règle générale : en chinois, le mot interrogatif occupe \emph{la place exacte du mot-réponse}, il ne se déplace jamais en tête de phrase.
\end{attention}

\begin{exemplebox}[Question — réponse : observer le mécanisme]
\zh{人们在哪儿买东西？}\\
\emph{Rénmen zài nǎr mǎi dōngxi ?}\\
« Où les gens achètent-ils des choses ? »\\[4pt]
\zh{人们在市场买东西。}\\
\emph{Rénmen zài shìchǎng mǎi dōngxi.}\\
« Les gens achètent des choses au marché. »\\[6pt]
Observez : seul le mot \zh{哪儿} a disparu, remplacé par \zh{市场}. Le reste de la phrase est identique. C'est la règle d'or de la réponse en chinois.
\end{exemplebox}

\courssub{Questions de compréhension sur le texte}

\begin{textebox}[Rappel : la phrase clé du texte]
\zh{经济的发展让人民的生活越来越好。}\\
\emph{Jīngjì de fāzhǎn ràng rénmín de shēnghuó yuèláiyuè hǎo.}\\
« Le développement économique améliore de plus en plus la vie du peuple. »\\[4pt]
Cette phrase utilise la structure \zh{让 + quelqu'un + verbe/adjectif} (faire que, permettre que) et la structure \zh{越来越 + adjectif} (de plus en plus...). Retenez-la : elle servira de modèle pour parler du Cameroun.
\end{textebox}

Répondez maintenant aux quatre questions en phrases complètes, en reprenant les mots de la question.

\begin{exoresolu}[Les quatre questions de compréhension]
\textbf{Énoncé.} Répondez en chinois, en phrases complètes :\\
1. \zh{中国的经济发展怎么样？} 2. \zh{人们在哪儿买东西？} 3. \zh{中国和哪些国家做生意？} 4. \zh{经济的发展让人民的生活怎么样？}
\tcblower
\textbf{Solution détaillée.}\\
1. \zh{中国的经济发展很快。} \emph{Zhōngguó de jīngjì fāzhǎn hěn kuài.} « L'économie chinoise se développe très vite. » — \zh{怎么样} est remplacé par l'adjectif \zh{很快}.\\
2. \zh{人们在市场买东西，也在网上买东西。} \emph{Rénmen zài shìchǎng mǎi dōngxi, yě zài wǎngshang mǎi dōngxi.} « Les gens achètent au marché, et aussi sur Internet. » — \zh{哪儿} est remplacé par \zh{市场} et \zh{网上}, précédés de \zh{在}.\\
3. \zh{中国和非洲国家做生意。} \emph{Zhōngguó hé Fēizhōu guójiā zuò shēngyi.} « La Chine fait du commerce avec les pays africains. » — \zh{哪些国家} est remplacé par \zh{非洲国家}.\\
4. \zh{经济的发展让人民的生活越来越好。} \emph{Jīngjì de fāzhǎn ràng rénmín de shēnghuó yuèláiyuè hǎo.} « Le développement économique rend la vie du peuple de plus en plus bonne. » — \zh{怎么样} est remplacé par \zh{越来越好}.
\end{exoresolu}

\begin{center}
\begin{tabularx}{\linewidth}{|X|X|X|}
\hline
\rowcolor{popblueL} Mot interrogatif & Question & Réponse (mot de remplacement) \\
\hline
\zh{怎么样} (comment ?) & \zh{经济发展怎么样？} & adjectif : \zh{很快} / \zh{越来越好} \\
\hline
\zh{哪儿} (où ?) & \zh{人们在哪儿买东西？} & \zh{在} + lieu : \zh{在市场} / \zh{在网上} \\
\hline
\zh{哪些} (lesquels ?) & \zh{和哪些国家做生意？} & nom + \zh{国家} : \zh{非洲国家} \\
\hline
\zh{谁} (qui ?) & \zh{谁去市场？} & personne : \zh{我妈妈去市场。} \\
\hline
\end{tabularx}
\end{center}

\begin{attention}[Ne pas oublier 在 devant le lieu]
Quand la réponse à \zh{哪儿} est un lieu, il faut garder la préposition \zh{在} (à, dans, sur) : on dit \zh{在市场买东西}, jamais \zh{市场买东西}. En français, « au marché » contient déjà la préposition ; en chinois, elle est séparée et obligatoire.
\end{attention}

\courssub{Carte mentale de compréhension du texte}

\begin{popfigure}
\begin{tikzpicture}[node distance=0.9cm and 1.6cm,
  boite/.style={draw=popblue, fill=popblueL, rounded corners=3pt, align=center, inner sep=5pt, font=\small},
  res/.style={draw=popgreen, fill=popgreenL, rounded corners=3pt, align=center, inner sep=5pt, font=\small}]
\node[boite, fill=poporangeL, draw=poporange, align=center] (t){Texte : l'économie\\\zh{经济}};
\node[boite, below left=of t, align=center] (qui){Qui ?\\\zh{中国 / 喀麦隆}};
\node[boite, below=of t, align=center] (quoi){Quoi ?\\\zh{经济发展很快}\\\zh{在市场、在网上买东西}};
\node[res, below right=of t, align=center] (r){Résultat ?\\\zh{生活越来越好}};
\draw[-{Stealth}, thick, popdark] (t) -- (qui);
\draw[-{Stealth}, thick, popdark] (t) -- (quoi);
\draw[-{Stealth}, thick, popdark] (t) -- (r);
\draw[-{Stealth}, thick, popgreen] (quoi) -- (r);
\end{tikzpicture}
\legende{Organigramme de compréhension : du texte vers les trois questions essentielles (qui ? quoi ? quel résultat ?).}
\end{popfigure}

\courssub{Questions ouvertes : et au Cameroun ?}

La langue sert à parler de \emph{son} monde. Transposons le texte au Cameroun avec la structure \zh{越来越} (de plus en plus), qui s'emploie devant un adjectif : \zh{越来越 + adjectif}.

\begin{propriete}[La structure 越来越 + adjectif]
\zh{越来越} (yuèláiyuè, « de plus en plus ») se place \textbf{directement devant l'adjectif}, \textbf{sans} \zh{很} : on dit \zh{越来越好} (de mieux en mieux), jamais \zh{越来越很好}. L'adjectif seul suffit, car \zh{越来越} exprime déjà l'intensité croissante.\\
Schéma : Sujet + \zh{越来越} + adjectif.\\
\zh{天气越来越热。} \emph{Tiānqì yuèláiyuè rè.} « Il fait de plus en plus chaud. »
\end{propriete}

\begin{exemplebox}[Parler de l'économie camerounaise]
\zh{喀麦隆的经济怎么样？}\\
\emph{Kāmàilóng de jīngjì zěnmeyàng ?}\\
« Comment va l'économie camerounaise ? »\\[4pt]
Réponses guidées possibles :\\
\zh{喀麦隆的经济发展越来越快。} \emph{Kāmàilóng de jīngjì fāzhǎn yuèláiyuè kuài.} « L'économie camerounaise se développe de plus en plus vite. »\\
\zh{喀麦隆和中国做生意。} \emph{Kāmàilóng hé Zhōngguó zuò shēngyi.} « Le Cameroun fait du commerce avec la Chine. »\\
\zh{在杜阿拉，人们在市场买东西。} \emph{Zài Dù'ālā, rénmen zài shìchǎng mǎi dōngxi.} « À Douala, les gens achètent au marché. »
\end{exemplebox}

\begin{experience}[Reformulation : résumer le texte en deux phrases]
Par groupes de deux, chaque élève résume le texte en \textbf{deux phrases} avec ses propres mots, sans regarder le texte. Modèle attendu :\\
1. Une phrase sur l'économie avec \zh{发展} ou \zh{越来越} : \zh{中国的经济发展很快。}\\
2. Une phrase sur la vie des gens avec \zh{让} ou \zh{越来越} : \zh{人民的生活越来越好。}\\
Puis chaque binôme adapte son résumé au Cameroun : \zh{喀麦隆的经济……} Le reste de la classe vérifie que les deux structures sont correctes.
\end{experience}

\begin{exercice}[À vous de répondre]
Répondez en chinois, en phrases complètes :\\
1. \zh{喀麦隆的经济发展怎么样？}\\
2. \zh{你在哪儿买东西？}\\
3. \zh{喀麦隆和哪些国家做生意？}\\
4. Traduisez en chinois : « La vie des gens est de plus en plus bonne. »
\end{exercice}

\begin{corrige}[Exercice : À vous de répondre]
1. \zh{喀麦隆的经济发展越来越快。} \emph{Kāmàilóng de jīngjì fāzhǎn yuèláiyuè kuài.} « L'économie camerounaise se développe de plus en plus vite. »\\
2. \zh{我在市场买东西。} \emph{Wǒ zài shìchǎng mǎi dōngxi.} « J'achète au marché. » (on accepte aussi \zh{在网上} : sur Internet).\\
3. \zh{喀麦隆和中国做生意。} \emph{Kāmàilóng hé Zhōngguó zuò shēngyi.} « Le Cameroun fait du commerce avec la Chine. »\\
4. \zh{人民的生活越来越好。} \emph{Rénmín de shēnghuó yuèláiyuè hǎo.} — Attention : \zh{越来越} se place juste devant l'adjectif \zh{好}, sans \zh{很}.
\end{corrige}

\begin{aretenir}[L'essentiel de la section]
\begin{itemize}
\item Pour répondre à une question en chinois : recopier la question, remplacer le mot interrogatif par la réponse, sans changer l'ordre des mots.
\item \zh{怎么样} et \zh{哪儿} restent dans la phrase, à la place de la réponse — jamais en tête comme « comment » et « où » en français.
\item Devant un lieu, garder \zh{在} : \zh{在市场}, \zh{在网上}.
\item La structure \zh{越来越 + adjectif} exprime « de plus en plus » : \zh{越来越好}, \zh{越来越快}. Ne pas ajouter \zh{很} devant l'adjectif.
\item La structure \zh{A 让 B + adjectif} exprime « A fait que B est... » : \zh{经济的发展让人民的生活越来越好。}
\end{itemize}
\end{aretenir}

\coursec{Production guidée et ancrage Cameroun-Chine}

Après avoir répondu aux questions de compréhension du texte support, vous savez maintenant \emph{lire} un texte chinois sur l'économie. Il reste l'étape la plus importante : \textbf{produire vous-même}. Dans cette dernière partie, vous allez passer du statut de lecteur à celui d'auteur : d'abord en imitant un modèle, puis en complétant une trame, enfin en écrivant librement quelques phrases sur l'économie du Cameroun. C'est exactement la démarche attendue à l'examen : comprendre, imiter, produire.

\courssub{Observer le modèle : un mini-paragraphe sur l'économie du Cameroun}

Lisons d'abord le paragraphe modèle. Observez-le avant toute explication : combien de phrases ? De quoi parle chaque phrase ?

\begin{textebox}[Modèle de paragraphe : l'économie du Cameroun]
喀麦隆的经济也在发展。\\
Kāmàilóng de jīngjì yě zài fāzhǎn.\\
L'économie du Cameroun se développe aussi.\\[4pt]
很多人在市场做生意。\\
Hěn duō rén zài shìchǎng zuò shēngyi.\\
Beaucoup de gens font du commerce au marché.\\[4pt]
中国和喀麦隆是好伙伴。\\
Zhōngguó hé Kāmàilóng shì hǎo huǒbàn.\\
La Chine et le Cameroun sont de bons partenaires.\\[4pt]
我希望喀麦隆的经济越来越好。\\
Wǒ xīwàng Kāmàilóng de jīngjì yuèláiyuè hǎo.\\
J'espère que l'économie du Cameroun ira de mieux en mieux.
\end{textebox}

\begin{definition}[Le vocabulaire clé du modèle]
\begin{center}
\begin{tabularx}{\linewidth}{|X|X|X|X|}
\hline
\rowcolor{popblueL} Caractères & Pinyin & Nature & Traduction \\
\hline
经济 & jīngjì & nom & économie \\
\hline
发展 & fāzhǎn & verbe / nom & (se) développer, développement \\
\hline
市场 & shìchǎng & nom & marché \\
\hline
做生意 & zuò shēngyi & expression verbale & faire du commerce, des affaires \\
\hline
伙伴 & huǒbàn & nom & partenaire, compagnon \\
\hline
希望 & xīwàng & verbe / nom & espérer, espoir \\
\hline
越来越好 & yuèláiyuè hǎo & structure & de mieux en mieux \\
\hline
\end{tabularx}
\end{center}
\end{definition}

\courssub{Analyser la structure du modèle : la chaîne en quatre étapes}

Regardez bien : ce paragraphe n'est pas écrit au hasard. Il suit une \cle{chaîne logique} très simple que vous pouvez réutiliser pour n'importe quel pays, n'importe quel thème économique.

\begin{methode}[Écrire un mini-paragraphe : la chaîne pays + économie + activités + souhait]
Pour rédiger un mini-paragraphe économique en chinois, suivez toujours ces quatre étapes :
\begin{enumerate}
\item \textbf{Le pays et son économie} : \zh{喀麦隆的经济也在发展。} — on annonce le sujet général (pays + \zh{的经济} + verbe).
\item \textbf{Les activités des gens} : \zh{很多人在市场做生意。} — on donne un exemple concret (\zh{很多人} + lieu + activité).
\item \textbf{Les partenaires} : \zh{中国和喀麦隆是好伙伴。} — on situe le pays dans le monde (A \zh{和} B \zh{是} + relation).
\item \textbf{Le souhait final} : \zh{我希望……越来越好。} — on conclut par un vœu personnel avec \zh{希望} et la structure \zh{越来越} + adjectif.
\end{enumerate}
Cette chaîne fonctionne aussi pour la ville, l'école, l'agriculture : il suffit de changer le sujet.
\end{methode}

\begin{propriete}[La structure 越来越 + adjectif : « de plus en plus »]
\textbf{Forme} : Sujet + \zh{越来越} + adjectif.\\
\textbf{Emploi} : exprime une évolution progressive, une gradation dans le temps.\\
\textbf{Exemples} :
\begin{itemize}
\item \zh{天气越来越热。} \emph{Tiānqì yuèláiyuè rè.} « Il fait de plus en plus chaud. »
\item \zh{他的中文越来越好。} \emph{Tā de Zhōngwén yuèláiyuè hǎo.} « Son chinois est de mieux en mieux. »
\item \zh{市场的人越来越多。} \emph{Shìchǎng de rén yuèláiyuè duō.} « Il y a de plus en plus de monde au marché. »
\end{itemize}
\textbf{Exception importante} : on ne dit jamais \zh{很} devant l'adjectif dans cette structure. On dit \zh{越来越好}, jamais \zh{越来越很好}. La gradation est déjà contenue dans \zh{越来越}.
\end{propriete}

\begin{attention}[希望 ne prend jamais de « que »]
En français, on dit « j'espère \textbf{que} l'économie ira mieux ». En chinois, \zh{希望} (xīwàng, espérer) est suivi \textbf{directement d'une phrase complète}, sans aucun mot de liaison :
\begin{itemize}
\item Correct : \zh{我希望喀麦隆的经济越来越好。} \emph{Wǒ xīwàng Kāmàilóng de jīngjì yuèláiyuè hǎo.}
\item Faux : ajouter un mot entre \zh{希望} et la phrase — il n'existe pas de « que » en chinois.
\end{itemize}
Autre piège de francophone : ne confondez pas \zh{希望} (espérer, souhait réaliste) avec \zh{想} (vouloir, avoir envie de). \zh{我想去中国。} « Je veux aller en Chine » ; \zh{我希望去中国。} « J'espère aller en Chine ».
\end{attention}

\courssub{Visualiser la chaîne économique}

Le paragraphe modèle repose sur une chaîne économique que l'on retrouve dans la réalité camerounaise : le marché nourrit les entreprises, les entreprises nourrissent le commerce, le commerce nourrit le développement, et le développement améliore la vie.

\begin{popfigure}
\begin{center}
\begin{tikzpicture}[node distance=0.55cm and 0.9cm]
\node[draw=popblue, fill=popblueL, rounded corners, thick, minimum width=1.9cm, minimum height=1cm, align=center] (a) {市场\\ \emph{\small shìchǎng}\\ \small marché};
\node[draw=poporange, fill=poporangeL, rounded corners, thick, minimum width=1.9cm, minimum height=1cm, align=center, right=of a] (b) {公司\\ \emph{\small gōngsī}\\ \small entreprise};
\node[draw=poppurple, fill=poppurpleL, rounded corners, thick, minimum width=1.9cm, minimum height=1cm, align=center, right=of b] (c) {贸易\\ \emph{\small màoyì}\\ \small commerce};
\node[draw=popgreen, fill=popgreenL, rounded corners, thick, minimum width=1.9cm, minimum height=1cm, align=center, right=of c] (d) {发展\\ \emph{\small fāzhǎn}\\ \small développement};
\node[draw=popgold, fill=popgoldL, rounded corners, thick, minimum width=2.4cm, minimum height=1cm, align=center, right=of d] (e) {生活越来越好\\ \emph{\small shēnghuó yuèláiyuè hǎo}\\ \small la vie s'améliore};
\draw[-{Stealth}, very thick, popdark] (a) -- (b);
\draw[-{Stealth}, very thick, popdark] (b) -- (c);
\draw[-{Stealth}, very thick, popdark] (c) -- (d);
\draw[-{Stealth}, very thick, popdark] (d) -- (e);
\end{tikzpicture}
\end{center}
\legende{La chaîne économique : du marché local à l'amélioration de la vie}
\end{popfigure}

\courssub{S'entraîner avec la trame à trous}

Avant d'écrire seul, complétez la trame. C'est un exercice de \cle{production guidée} : la structure est donnée, vous fournissez le contenu.

\begin{exoresolu}[Compléter la trame à trous]
\textbf{Énoncé.} Complétez la trame suivante avec les mots de votre choix (plusieurs réponses possibles) :\\
\zh{\_\_\_的经济\_\_\_。很多人在\_\_\_买东西。\_\_\_和\_\_\_是伙伴。我希望\_\_\_越来越好。}
\tcblower
\textbf{Solution détaillée.} Une complétion possible :\\
\zh{喀麦隆的经济在发展。很多人在超市买东西。中国和喀麦隆是伙伴。我希望我们的生活越来越好。}\\
\emph{Kāmàilóng de jīngjì zài fāzhǎn. Hěn duō rén zài chāoshì mǎi dōngxi. Zhōngguó hé Kāmàilóng shì huǒbàn. Wǒ xīwàng wǒmen de shēnghuó yuèláiyuè hǎo.}\\
« L'économie du Cameroun se développe. Beaucoup de gens achètent des choses au supermarché. La Chine et le Cameroun sont partenaires. J'espère que notre vie ira de mieux en mieux. »\\
\textbf{Méthode} : 1) repérer le rôle de chaque trou (sujet, verbe, lieu, pays) ; 2) choisir un mot du glossaire de la leçon ; 3) vérifier que chaque phrase a un sujet et un verbe ; 4) relire à voix haute en faisant attention aux tons.
\end{exoresolu}

Voici un tableau des mots que vous pouvez glisser dans les trous, classés par fonction :

\begin{center}
\begin{tabularx}{\linewidth}{|X|X|X|X|}
\hline
\rowcolor{popblueL} Fonction dans la phrase & Caractères & Pinyin & Traduction \\
\hline
Pays / lieu & 喀麦隆 / 中国 / 杜阿拉 & Kāmàilóng / Zhōngguó / Dùālā & Cameroun / Chine / Douala \\
\hline
Verbe d'évolution & 在发展 / 越来越好 & zài fāzhǎn / yuèláiyuè hǎo & se développe / de mieux en mieux \\
\hline
Lieu d'activité & 市场 / 超市 / 商店 & shìchǎng / chāoshì / shāngdiàn & marché / supermarché / magasin \\
\hline
Activité & 买东西 / 做生意 / 工作 & mǎi dōngxi / zuò shēngyi / gōngzuò & acheter / commercer / travailler \\
\hline
Relation & 伙伴 / 朋友 & huǒbàn / péngyou & partenaire / ami \\
\hline
Souhait & 生活 / 经济 / 农业 & shēnghuó / jīngjì / nóngyè & la vie / l'économie / l'agriculture \\
\hline
\end{tabularx}
\end{center}

\courssub{Texte d'entraînement : au marché de Yaoundé}

Lisons maintenant un court texte original qui réutilise tout le vocabulaire de la leçon dans une situation camerounaise authentique. Lisez-le deux fois : une fois pour le sens global, une fois en soulignant les mots du glossaire.

\begin{textebox}[雅温得的莫克洛市场 — Le marché Mokolo de Yaoundé]
星期六早上，雅温得的莫克洛市场人很多。\\
Xīngqīliù zǎoshang, Yǎwēndé de Mòkèluò shìchǎng rén hěn duō.\\
Samedi matin, il y a beaucoup de monde au marché Mokolo de Yaoundé.\\[4pt]
很多人在这里做生意：有的人卖水果，有的人卖鱼，还有的人卖衣服。\\
Hěn duō rén zài zhèlǐ zuò shēngyi: yǒu de rén mài shuǐguǒ, yǒu de rén mài yú, hái yǒu de rén mài yīfu.\\
Beaucoup de gens y font du commerce : certains vendent des fruits, d'autres vendent du poisson, d'autres encore vendent des vêtements.\\[4pt]
市场里也有中国商品：手机、鞋和电视机。\\
Shìchǎng lǐ yě yǒu Zhōngguó shāngpǐn: shǒujī, xié hé diànshìjī.\\
Dans le marché, il y a aussi des produits chinois : téléphones portables, chaussures et téléviseurs.\\[4pt]
中国和喀麦隆的贸易越来越多，两个国家是好伙伴。\\
Zhōngguó hé Kāmàilóng de màoyì yuèláiyuè duō, liǎng ge guójiā shì hǎo huǒbàn.\\
Le commerce entre la Chine et le Cameroun est de plus en plus important, les deux pays sont de bons partenaires.\\[4pt]
我希望喀麦隆的经济越来越好，人们的生活也越来越好。\\
Wǒ xīwàng Kāmàilóng de jīngjì yuèláiyuè hǎo, rénmen de shēnghuó yě yuèláiyuè hǎo.\\
J'espère que l'économie du Cameroun ira de mieux en mieux, et que la vie des gens s'améliorera aussi.
\end{textebox}

\textbf{Glossaire du texte.}

\begin{center}
\begin{tabularx}{\linewidth}{|X|X|X|X|}
\hline
\rowcolor{popblueL} Caractères & Pinyin & Nature & Traduction \\
\hline
星期六 & xīngqīliù & nom & samedi \\
\hline
卖 & mài & verbe & vendre \\
\hline
水果 & shuǐguǒ & nom & fruits \\
\hline
衣服 & yīfu & nom & vêtements \\
\hline
商品 & shāngpǐn & nom & marchandise, produit \\
\hline
手机 & shǒujī & nom & téléphone portable \\
\hline
贸易 & màoyì & nom & commerce (international) \\
\hline
国家 & guójiā & nom & pays, État \\
\hline
人们 & rénmen & nom & les gens \\
\hline
\end{tabularx}
\end{center}

\textbf{Questions rapides.} 1) Quel jour le marché est-il décrit ? 2) Citez deux produits vendus au marché. 3) Quels produits chinois voit-on au marché ? 4) Quel est le souhait du narrateur ?

\courssub{Production individuelle : à vous d'écrire}

\begin{exercice}[Mon mini-paragraphe sur l'économie du Cameroun]
Rédigez un paragraphe de \textbf{4 à 5 phrases} sur l'économie du Cameroun, en suivant la chaîne : pays + économie + activités des gens + partenaire + souhait.\\
Contraintes :
\begin{itemize}
\item utilisez au moins une fois la structure \zh{越来越} + adjectif ;
\item utilisez au moins une fois \zh{希望} suivi d'une phrase complète ;
\item mentionnez un marché local (Mokolo, le marché de Douala, celui de votre ville), l'agriculture (cacao, café, bananes) ou le commerce avec la Chine ;
\item écrivez les caractères appris dans la leçon, avec le pinyin au-dessus si besoin.
\end{itemize}
\textbf{Amorce possible} : \zh{喀麦隆的经济……} / \zh{很多人在……} / \zh{中国和喀麦隆……} / \zh{我希望……}
\end{exercice}

\begin{corrige}[Exemple de production attendue]
\zh{喀麦隆的经济在发展。很多人在莫克洛市场做生意。农民种可可和咖啡。中国和喀麦隆是好伙伴，两国的贸易越来越多。我希望喀麦隆人的生活越来越好。}\\
\emph{Kāmàilóng de jīngjì zài fāzhǎn. Hěn duō rén zài Mòkèluò shìchǎng zuò shēngyi. Nóngmín zhòng kěkě hé kāfēi. Zhōngguó hé Kāmàilóng shì hǎo huǒbàn, liǎng guó de màoyì yuèláiyuè duō. Wǒ xīwàng Kāmàilóng rén de shēnghuó yuèláiyuè hǎo.}\\
« L'économie du Cameroun se développe. Beaucoup de gens font du commerce au marché Mokolo. Les paysans cultivent le cacao et le café. La Chine et le Cameroun sont de bons partenaires, le commerce entre les deux pays est de plus en plus important. J'espère que la vie des Camerounais ira de mieux en mieux. »\\
\textbf{Points de vigilance du correcteur} : ordre des mots (sujet + lieu avec \zh{在} + verbe), absence de « que » après \zh{希望}, pas de \zh{很} dans \zh{越来越好}, tons corrects sur \zh{mǎi} (acheter, 3\up{e} ton) et \zh{mài} (vendre, 4\up{e} ton).
\end{corrige}

\begin{attention}[买 et 卖 : le piège classique des francophones]
\zh{买} (mǎi, 3\up{e} ton) = \textbf{acheter} ; \zh{卖} (mài, 4\up{e} ton) = \textbf{vendre}. Mêmes syllabes, tons différents, sens opposés ! Moyen mnémotechnique : \zh{卖} a un petit \zh{十} en haut, comme une étiquette de prix posée sur la marchandise — celui qui \emph{vend} met une étiquette. À l'oral, un mauvais ton transforme « j'achète du poisson » en « je vends du poisson » : entraînez-vous avec la paire \zh{我买东西。} / \zh{他卖东西。}
\end{attention}

\courssub{Échange en binômes : lecture mutuelle et correction}

\begin{experience}[Binômes : je lis, tu corriges]
Travaillez par deux, en trois étapes :
\begin{enumerate}
\item \textbf{Lecture mutuelle} : chaque élève lit son paragraphe à voix haute à son partenaire, lentement, en soignant les tons.
\item \textbf{Correction ciblée} : le partenaire écoute et note uniquement deux types d'erreurs : (a) les tons faux (surtout \zh{mǎi}/\zh{mài}, \zh{shìchǎng}, \zh{yuèláiyuè}) ; (b) les caractères mal écrits (vérifier \zh{经}, \zh{济}, \zh{商}, \zh{越}).
\item \textbf{Relecture finale} : chacun réécrit la phrase corrigée et la relit. Le binôme choisit ensuite la meilleure phrase des deux et la présente à la classe.
\end{enumerate}
Règle d'or de la correction entre pairs : on signale l'erreur, on ne la corrige pas à la place de l'autre — c'est à l'auteur de trouver la bonne forme.
\end{experience}

\courssub{Complément examen : présenter un texte en trois phrases}

\begin{methode}[Présenter un texte en 3 phrases — compétence d'épreuve de compréhension écrite]
À l'examen, on peut vous demander de résumer ou de présenter un texte. Retenez ce plan en trois phrases, valable pour tout texte :
\begin{enumerate}
\item \textbf{Le thème} : « Ce texte parle de… » — \zh{这个文本讲……} \emph{Zhège wénběn jiǎng…} Exemple : \zh{这个文本讲喀麦隆的市场。}
\item \textbf{Les idées principales} : une ou deux informations essentielles — \zh{很多人在市场做生意，市场里也有中国商品。}
\item \textbf{La conclusion} : l'idée finale ou le message — \zh{作者希望喀麦隆的经济越来越好。}
\end{enumerate}
Entraînez-vous à appliquer ce plan au texte « Le marché Mokolo » ci-dessus : c'est exactement le type de réponse courte et structurée que le correcteur attend.
\end{methode}

\begin{savaistu}[Le port de Kribi, symbole de la coopération sino-camerounaise]
Le port en eau profonde de Kribi, sur la côte atlantique du Cameroun, a été construit avec la participation d'entreprises chinoises. Il est aujourd'hui l'un des symboles les plus visibles de la coopération économique entre la Chine et le Cameroun : par lui transitent des marchandises qui arrivent ensuite sur les marchés de tout le pays, comme celles que l'on voit au marché Mokolo. En chinois, on dit \zh{合作} (hézuò, coopération) : \zh{中国和喀麦隆合作建设克里比港口。} « La Chine et le Cameroun coopèrent pour construire le port de Kribi. »
\end{savaistu}

\begin{aretenir}[L'essentiel de la section]
\begin{itemize}
\item Un mini-paragraphe économique suit la chaîne : \textbf{pays + économie + activités des gens + partenaire + souhait}.
\item \zh{越来越} + adjectif = « de plus en plus », jamais avec \zh{很}.
\item \zh{希望} est suivi directement d'une phrase complète, sans « que ».
\item \zh{买} (mǎi, acheter) et \zh{卖} (mài, vendre) se distinguent par le ton.
\item À l'examen, présenter un texte en 3 phrases : thème, idées principales, conclusion.
\end{itemize}
\end{aretenir}

\newpage

\coursec{Activité d'intégration}

\courssub{Objectifs de l'activité}

À l'issue de cette activité d'intégration, l'élève sera capable de :
\begin{itemize}
\item mobiliser dans une production orale authentique au moins six mots du glossaire de la leçon (\zh{经济} \emph{jīngjì} « économie », \zh{市场} \emph{shìchǎng} « marché », \zh{公司} \emph{gōngsī} « entreprise », \zh{贸易} \emph{màoyì} « commerce », \zh{发展} \emph{fāzhǎn} « développer », \zh{工作} \emph{gōngzuò} « travail », etc.) ;
\item réutiliser la trame du texte support (présentation d'un lieu, d'une activité, d'un échange) en l'adaptant au contexte camerounais ;
\item prononcer correctement les tons et le rythme de phrases courtes de niveau HSK 2-3 ;
\item écrire au tableau un caractère du thème en respectant l'ordre des traits (haut → bas, gauche → droite, horizontal avant vertical) ;
\item écouter les exposés des camarades, poser une question simple et voter de façon argumentée.
\end{itemize}

\courssub{Situation}

Dans le cadre de la semaine « Cameroun-Chine » organisée par ton lycée de Yaoundé, le club de chinois prépare un petit « marché de l'économie » : chaque binôme présente en deux minutes un mini-exposé sur l'économie de notre pays devant les élèves de Seconde et un professeur invité. Une affiche annonce l'événement.

\begin{textebox}[L'affiche du club de chinois]
中文俱乐部欢迎你！\\
Zhōngwén jùlèbù huānyíng nǐ !\\
« Le club de chinois te souhaite la bienvenue ! »\\[4pt]
主题：我们国家的经济\\
Zhǔtí : Wǒmen guójiā de jīngjì.\\
« Thème : l'économie de notre pays. »\\[4pt]
说一说喀麦隆的市场、公司和贸易！\\
Shuō yi shuō Kāmàilóng de shìchǎng, gōngsī hé màoyì !\\
« Parlez des marchés, des entreprises et du commerce du Cameroun ! »\\[4pt]
时间：星期五下午三点。地点：二号教室。\\
Shíjiān : xīngqīwǔ xiàwǔ sān diǎn. Dìdiǎn : èr hào jiàoshì.\\
« Horaire : vendredi à 15 heures. Lieu : salle n° 2. »
\end{textebox}

Ton binôme est inscrit. Tu dois préparer un exposé de cinq phrases intitulé \zh{我们国家的经济} (\emph{Wǒmen guójiā de jīngjì}, « L'économie de notre pays »), en t'inspirant du texte support de la leçon, mais en parlant du Cameroun : ses marchés (Mokolo, Mvog-Mbi, le marché central de Douala), ses entreprises, son commerce avec la Chine.

\courssub{Consignes}

\begin{enumerate}
\item \textbf{Relecture ciblée du texte support.} Relis le texte de la leçon et souligne : (a) trois mots du glossaire que tu veux réutiliser ; (b) une structure de phrase que tu peux copier en changeant seulement un ou deux mots (par exemple \zh{……有很多……} « il y a beaucoup de … »).
\item \textbf{Choix du lexique.} Dans le tableau de vocabulaire, choisis au moins \cle{six mots} du thème « économie » et écris-les sur ta fiche avec le pinyin et les tons. Vérifie la prononciation avec ton binôme.
\item \textbf{Rédaction de la trame.} Rédige cinq phrases en suivant ce plan : phrase 1 = présentation générale (\zh{喀麦隆的经济……}) ; phrase 2 = les marchés ; phrase 3 = les entreprises et le travail ; phrase 4 = le commerce avec la Chine ; phrase 5 = une conclusion avec \zh{发展} ou \zh{越来越……}. Chaque phrase doit rester courte (8 à 15 caractères).
\item \textbf{Relecture grammaticale.} Vérifie : l'ordre Sujet + Verbe + Complément ; la place du complément de lieu avant le verbe (\zh{在市场上} « au marché ») ; l'emploi de \zh{的} entre le nom et son complément du nom (\zh{我们国家的经济}).
\item \textbf{Préparation de l'écriture au tableau.} Choisis un caractère du glossaire (par exemple \zh{市} ou \zh{贸}) et entraîne-toi à le tracer en respectant l'ordre des traits : pour \zh{市}, le point en haut, puis le trait horizontal, ensuite le vertical et le crochet ; pour \zh{贸}, la partie haute avant la clé \zh{贝} (la coquille, clé de la valeur et de l'argent) en bas.
\item \textbf{Présentation orale.} Présentez l'exposé à deux voix (alternez les phrases), à voix haute, en soignant les tons. Un des deux élèves écrit le caractère au tableau pendant que l'autre conclut.
\item \textbf{Écoute active et vote.} Pendant les exposés des autres binômes, note un mot du glossaire que tu as reconnu et prépare une question simple (\zh{这个市场在哪里？} « Où est ce marché ? »). À la fin, la classe vote pour l'exposé le plus clair : justifie ton vote en une phrase en français.
\end{enumerate}

\courssub{Ressources à mobiliser}

\begin{aretenir}[Boîte à outils de l'exposé]
\begin{itemize}
\item \textbf{Lexique du thème :} \zh{经济} \emph{jīngjì} « économie » ; \zh{市场} \emph{shìchǎng} « marché » ; \zh{公司} \emph{gōngsī} « entreprise » ; \zh{贸易} \emph{màoyì} « commerce » ; \zh{发展} \emph{fāzhǎn} « (se) développer » ; \zh{工作} \emph{gōngzuò} « travail, travailler » ; \zh{买} \emph{mǎi} « acheter » ; \zh{卖} \emph{mài} « vendre » ; \zh{贵/便宜} \emph{guì/piányi} « cher/bon marché ».
\item \textbf{Structures utiles :} Sujet + \zh{有} + nom (« il y a ») : \zh{喀麦隆有很多大市场。} ; \zh{越来越} + adjectif (« de plus en plus ») : \zh{越来越重要} « de plus en plus important » ; complément de lieu avant le verbe : \zh{人们在市场买东西。} « Les gens font leurs achats au marché. »
\item \textbf{Liaison :} \zh{和} \emph{hé} « et » (entre noms), \zh{也} \emph{yě} « aussi », \zh{因为……所以……} « parce que … donc … ».
\item \textbf{Écriture :} clé \zh{贝} (argent, valeur) dans \zh{贸}, \zh{贵} ; clé \zh{巾} (tissu) dans \zh{市} ; ordre des traits : haut → bas, gauche → droite, l'extérieur avant l'intérieur.
\end{itemize}
\end{aretenir}

\begin{attention}[Pièges fréquents des francophones]
\begin{itemize}
\item Ne dis pas \zh{经济是很好} : en chinois, l'adjectif seul fonctionne comme prédicat, sans \zh{是}. Dis \zh{经济很好} (\emph{jīngjì hěn hǎo}, « l'économie va bien »).
\item Attention à la paire \zh{买} \emph{mǎi} (3\up{e} ton, « acheter ») / \zh{卖} \emph{mài} (4\up{e} ton, « vendre ») : seul le ton les distingue.
\item \zh{和} ne relie pas deux phrases : pour « et » entre deux idées, utilise \zh{也} ou fais deux phrases.
\item Le complément de lieu se place avant le verbe : \zh{在市场买东西}, jamais \zh{买东西在市场}.
\end{itemize}
\end{attention}

\courssub{Proposition de production et corrigé commenté}

\begin{exoresolu}[Modèle d'exposé : 我们国家的经济]
Produis un exposé de cinq phrases avec au moins six mots du glossaire, puis présente-le à la classe.
\tcblower
\textbf{Production modèle :}\\[4pt]
\zh{我们国家的经济越来越大。}\\
\emph{Wǒmen guójiā de jīngjì yuè lái yuè dà.} « L'économie de notre pays devient de plus en plus importante. »\\
\zh{雅温得和杜阿拉有很多大市场，人们在那里买东西。}\\
\emph{Yǎwēndé hé Dù'ālā yǒu hěn duō dà shìchǎng, rénmen zài nàli mǎi dōngxi.} « Yaoundé et Douala ont beaucoup de grands marchés, les gens y font leurs achats. »\\
\zh{很多公司在城市工作，也给年轻人工作。}\\
\emph{Hěn duō gōngsī zài chéngshì gōngzuò, yě gěi niánqīngrén gōngzuò.} « Beaucoup d'entreprises travaillent en ville et donnent aussi du travail aux jeunes. »\\
\zh{喀麦隆和中国有很多贸易。}\\
\emph{Kāmàilóng hé Zhōngguó yǒu hěn duō màoyì.} « Le Cameroun et la Chine ont beaucoup d'échanges commerciaux. »\\
\zh{我们希望国家的经济发展得越来越好。}\\
\emph{Wǒmen xīwàng guójiā de jīngjì fāzhǎn de yuè lái yuè hǎo.} « Nous espérons que l'économie du pays se développera de mieux en mieux. »\\[6pt]
\textbf{Commentaire des choix de langue :}
\begin{itemize}
\item Six mots du glossaire sont employés : \zh{经济}, \zh{市场}, \zh{公司}, \zh{工作}, \zh{贸易}, \zh{发展} (plus \zh{买}).
\item La phrase 1 reprend la structure \zh{越来越} + adjectif vue en leçon ; la phrase 5 montre l'emploi de \zh{得} après le verbe \zh{发展} pour introduire le degré (\zh{发展得越来越好}).
\item La phrase 2 place correctement le complément de lieu \zh{在那里} avant le verbe \zh{买}.
\item La phrase 4 relie deux noms propres avec \zh{和}, ce qui est correct ; les phrases restent courtes, donc faciles à prononcer avec des tons nets.
\item Caractère tracé au tableau : \zh{贸} — partie haute d'abord, puis la clé \zh{贝} en bas (horizontal, vertical, puis les deux derniers traits).
\end{itemize}
\end{exoresolu}

\courssub{Bilan de l'activité}

Cette activité a permis de réunir tous les savoir-faire de la leçon dans une tâche de communication réelle : comprendre et réutiliser la trame d'un texte, choisir et prononcer le vocabulaire de l'économie, respecter les structures de base de la phrase chinoise et écrire un caractère dans les règles. L'ancrage Cameroun-Chine (marchés de Yaoundé et de Douala, échanges commerciaux) donne du sens à l'apprentissage : le chinois devient un outil pour parler de son propre pays. Le vote final, justifié en français, développe enfin l'écoute active et l'esprit critique. Pour aller plus loin, les meilleurs exposés peuvent être affichés dans la salle de chinois ou enregistrés pour la journée portes ouvertes du lycée.

\newpage

\coursec{Méthodes et exercices résolus}

\courssub{Savoir-faire 1 : Lire et comprendre un court texte chinois sur l'économie}

\begin{methode}[Comprendre un texte économique en quatre passages]
\begin{enumerate}
\item \textbf{Lecture silencieuse globale.} Lis le texte une première fois sans dictionnaire. Repère le \cle{thème} (ici : l'économie, le commerce, le travail) et les mots déjà connus : \zh{经济} (jīngjì, économie), \zh{市场} (shìchǎng, marché), \zh{工作} (gōngzuò, travail).
\item \textbf{Repérage des mots-clés du champ lexical.} Souligne les noms (lieux, métiers, produits), les verbes d'action (\zh{买} mǎi, acheter ; \zh{卖} mài, vendre ; \zh{发展} fāzhǎn, se développer) et les chiffres ou quantités.
\item \textbf{Découpage des phrases.} En chinois, la phrase de base suit l'ordre \cle{Sujet + Verbe + Complément}. Les compléments de temps et de lieu se placent \textbf{avant} le verbe : \zh{喀麦隆的经济这几年发展得很快。} (Kāmàilóng de jīngjì zhè jǐ nián fāzhǎn de hěn kuài.) « L'économie du Cameroun s'est développée rapidement ces dernières années. »
\item \textbf{Vérification par la traduction mentale.} Reformule chaque phrase en français simple, puis réponds aux questions en t'appuyant \emph{uniquement} sur le texte, jamais sur tes connaissances personnelles.
\end{enumerate}
\end{methode}

\begin{textebox}[在杜阿拉市场]
莉莉的妈妈在杜阿拉的一个大市场卖水果。\\
Lìli de māma zài Dù'ālā de yí ge dà shìchǎng mài shuǐguǒ.\\
« La maman de Lili vend des fruits dans un grand marché de Douala. »\\
她的水果又新鲜又便宜，所以客人很多。\\
Tā de shuǐguǒ yòu xīnxiān yòu piányi, suǒyǐ kèrén hěn duō.\\
« Ses fruits sont à la fois frais et bon marché, donc il y a beaucoup de clients. »\\
生意越来越好，她很高兴。\\
Shēngyi yuè lái yuè hǎo, tā hěn gāoxìng.\\
« Les affaires vont de mieux en mieux, elle est très contente. »
\end{textebox}

\begin{exoresolu}[Compréhension globale d'un texte sur le marché]
Énoncé. Lis le texte « 在杜阿拉市场 » ci-dessus, puis réponds en français : (1) Où se passe la scène ? (2) Que vend la maman de Lìli ? (3) Pourquoi les affaires vont-elles bien ?
\tcblower
\textbf{Solution détaillée.}
\begin{enumerate}
\item \textbf{Question 1.} Le texte dit \zh{在杜阿拉的一个大市场} : la préposition \zh{在} (zài) + lieu indique où se déroule l'action. Réponse : la scène se passe \textbf{dans un grand marché de Douala}, au Cameroun.
\item \textbf{Question 2.} Le verbe \zh{卖} (mài, vendre) a pour complément \zh{水果} (shuǐguǒ, fruits). Réponse : elle \textbf{vend des fruits}.
\item \textbf{Question 3.} La conséquence est introduite par \zh{所以} (suǒyǐ, donc) : la cause précède, c'est \zh{水果又新鲜又便宜}. La structure \zh{又……又……} (yòu... yòu...) signifie « à la fois... et... ». Réponse : les affaires vont bien parce que ses fruits sont \textbf{à la fois frais et bon marché}, donc les clients sont nombreux.
\end{enumerate}
\textbf{Conclusion.} Pour répondre, on repère d'abord le repère grammatical correspondant à la question (où → \zh{在} + lieu ; quoi → complément du verbe ; pourquoi → \zh{因为/所以}), puis on cite la phrase exacte du texte.
\end{exoresolu}

\courssub{Savoir-faire 2 : Reconnaître, prononcer et employer le vocabulaire de l'économie}

\begin{definition}[Glossaire du thème : l'économie]
\begin{center}
\begin{tabularx}{\linewidth}{|X|X|X|X|}
\hline
\rowcolor{popblueL} Caractères & Pinyin & Nature & Traduction \\
\hline
经济 & jīngjì & nom & économie \\
\hline
发展 & fāzhǎn & verbe & se développer \\
\hline
市场 & shìchǎng & nom & marché \\
\hline
商店 & shāngdiàn & nom & magasin, boutique \\
\hline
公司 & gōngsī & nom & société, entreprise \\
\hline
生意 & shēngyi & nom & affaires, commerce \\
\hline
买 & mǎi & verbe & acheter \\
\hline
卖 & mài & verbe & vendre \\
\hline
价格 & jiàgé & nom & prix \\
\hline
贵 & guì & adjectif & cher \\
\hline
便宜 & piányi & adjectif & bon marché \\
\hline
钱 & qián & nom & argent \\
\hline
工作 & gōngzuò & nom / verbe & travail ; travailler \\
\hline
工人 & gōngrén & nom & ouvrier \\
\hline
老板 & lǎobǎn & nom & patron, chef d'entreprise \\
\hline
产品 & chǎnpǐn & nom & produit \\
\hline
客人 & kèrén & nom & client, invité \\
\hline
国家 & guójiā & nom & pays, État \\
\hline
\end{tabularx}
\end{center}
\end{definition}

\begin{methode}[Mémoriser et réemployer le lexique économique]
\begin{enumerate}
\item \textbf{Apprends les mots par familles.} Regroupe : l'argent (\zh{钱} qián, \zh{价格} jiàgé prix, \zh{贵} guì cher, \zh{便宜} piányi bon marché), le commerce (\zh{买} mǎi acheter, \zh{卖} mài vendre, \zh{商店} shāngdiàn magasin, \zh{公司} gōngsī société), le travail (\zh{工作} gōngzuò, \zh{工人} gōngrén ouvrier, \zh{老板} lǎobǎn patron).
\item \textbf{Fixe les tons par paires contrastées.} \zh{买} mǎi (3\textsuperscript{e} ton, acheter) et \zh{卖} mài (4\textsuperscript{e} ton, vendre) ne se distinguent que par le ton et le caractère : répète-les toujours ensemble.
\item \textbf{Emploie chaque mot dans une phrase modèle.} Exemple : \zh{这个东西太贵了，我不买。} (Zhège dōngxi tài guì le, wǒ bù mǎi.) « Cet objet est trop cher, je ne l'achète pas. »
\item \textbf{Utilise les classificateurs corrects.} Un magasin : \zh{一家商店} (yì jiā shāngdiàn) ; une société : \zh{一家公司} (yì jiā gōngsī) ; un objet : \zh{一个东西} (yí ge dōngxi).
\end{enumerate}
\end{methode}

\begin{exoresolu}[Vocabulaire : compléter un dialogue au marché]
Énoncé. Complète le dialogue avec les mots suivants : \zh{便宜} (piányi), \zh{卖} (mài), \zh{价格} (jiàgé), \zh{公司} (gōngsī).

A：这个市场\_\_\_\_什么？\\
B：这里\_\_\_\_水果、蔬菜和鱼。\\
A：这些东西的\_\_\_\_怎么样？\\
B：不贵，很\_\_\_\_。我姐姐的\_\_\_\_也常来这里买东西。
\tcblower
\textbf{Solution détaillée.}
\begin{enumerate}
\item Première phrase : après le sujet \zh{这个市场} il faut un verbe ; la réponse parle de ce qu'on vend → \zh{卖}. « Que vend ce marché ? »
\item Deuxième phrase : sujet \zh{这里} + verbe → \zh{卖}. Le verbe \zh{卖} se répète naturellement dans la question et la réponse en chinois.
\item Troisième phrase : \zh{……的\_\_\_\_怎么样} demande une qualité ; le sujet doit être un nom → \zh{价格} (prix). « Comment sont les prix de ces choses ? »
\item Quatrième phrase : après \zh{很} il faut un adjectif → \zh{便宜} (bon marché), cohérent avec \zh{不贵}.
\item Cinquième phrase : \zh{我姐姐的} + nom de lieu de travail → \zh{公司}. « La société de ma sœur aînée vient aussi souvent acheter ici. »
\end{enumerate}
\textbf{Dialogue complet :} \zh{A：这个市场卖什么？B：这里卖水果、蔬菜和鱼。A：这些东西的价格怎么样？B：不贵，很便宜。我姐姐的公司也常来这里买东西。}\\
(A : Zhège shìchǎng mài shénme ? B : Zhèlǐ mài shuǐguǒ, shūcài hé yú. A : Zhèxiē dōngxi de jiàgé zěnmeyàng ? B : Bú guì, hěn piányi. Wǒ jiějie de gōngsī yě cháng lái zhèlǐ mǎi dōngxi.)\\
« A : Que vend ce marché ? B : Ici on vend des fruits, des légumes et du poisson. A : Comment sont les prix ? B : Pas chers, très bon marché. La société de ma sœur aînée vient aussi souvent acheter ici. »
\end{exoresolu}

\courssub{Savoir-faire 3 : Identifier les clés (radicaux) et l'ordre des traits}

\begin{methode}[Analyser un caractère du thème économique]
\begin{enumerate}
\item \textbf{Repère la clé (radical).} Elle donne souvent le sens : la clé \zh{贝} (bèi, coquillage → argent, car les coquillages servaient de monnaie dans la Chine ancienne) se trouve dans \zh{贵} (cher), \zh{货} (marchandise), \zh{买} et \zh{卖} ; la clé \zh{亻} (personne, variante de \zh{人}) dans \zh{价} (prix), \zh{作} (dans \zh{工作}).
\item \textbf{Décompose le caractère} en parties : clé + élément phonétique. Exemple : \zh{价} = \zh{亻} (personne) + \zh{介} (jiè, indice de prononciation).
\item \textbf{Respecte l'ordre des traits} : de haut en bas, de gauche à droite, l'horizontal avant le vertical, l'extérieur avant l'intérieur, on ferme la boîte en dernier.
\item \textbf{Écris lentement en comptant les traits} et en prononçant le pinyin avec le ton à chaque copie.
\end{enumerate}
\end{methode}

\begin{exoresolu}[Clés et traits : 贵 et 卖]
Énoncé. Pour les caractères \zh{贵} (guì, cher) et \zh{卖} (mài, vendre) : (1) donne la clé ; (2) décris l'ordre d'écriture ; (3) écris une phrase avec chacun.
\tcblower
\textbf{Solution détaillée.}
\begin{enumerate}
\item \textbf{贵} (guì, 9 traits) : la clé est \zh{贝} (coquillage, symbole de l'argent), placée en bas. On écrit d'abord la partie supérieure de haut en bas (la barre verticale centrale, puis le cadre, puis la longue horizontale), puis la clé \zh{贝} en bas : le cadre (vertical puis horizontal-crochet), les deux traits intérieurs, et le point final en bas à droite. Règles appliquées : haut → bas, extérieur → intérieur.
\item \textbf{卖} (mài, 8 traits) : la clé est aussi \zh{贝}. On écrit d'abord la partie haute : l'horizontale, puis la verticale de \zh{十} (l'horizontal avant le vertical), puis le crochet et les deux points de la partie \zh{买}, et enfin \zh{贝} en bas. Ne pas confondre avec \zh{买} (mǎi, acheter), qui n'a pas le \zh{十} au-dessus.
\item Phrases : \zh{这个手机太贵了。} (Zhège shǒujī tài guì le.) « Ce téléphone est trop cher. » \zh{她在市场卖鱼。} (Tā zài shìchǎng mài yú.) « Elle vend du poisson au marché. »
\end{enumerate}
\textbf{Conclusion.} Retenir la clé \zh{贝} = famille de l'argent permet de deviner le sens de nombreux caractères économiques et d'éviter la confusion \zh{买}/\zh{卖}.
\end{exoresolu}

\courssub{Savoir-faire 4 : Répondre à des questions de compréhension en chinois}

\begin{methode}[Rédiger une réponse en chinois correct]
\begin{enumerate}
\item \textbf{Repère le mot interrogatif} : \zh{谁} (shéi, qui), \zh{什么} (shénme, quoi), \zh{哪儿/哪里} (nǎr/nǎli, où), \zh{为什么} (wèishénme, pourquoi), \zh{怎么样} (zěnmeyàng, comment).
\item \textbf{Recopie la structure de la question} en remplaçant le mot interrogatif par l'information du texte : \zh{她为什么很高兴？} → \zh{她很高兴，因为……}
\item \textbf{Respecte l'ordre des mots} : sujet + (complément de temps/lieu) + verbe + complément. Jamais de verbe conjugué, jamais d'ordre français.
\item \textbf{Relis} : vérifie les tons du pinyin, la ponctuation chinoise （。 ？） et la présence du classificateur si un nom est compté.
\end{enumerate}
\end{methode}

\begin{textebox}[中国经济和喀麦隆]
中国是一个经济发展很快的国家。\\
Zhōngguó shì yí ge jīngjì fāzhǎn hěn kuài de guójiā.\\
« La Chine est un pays dont l'économie se développe très vite. »\\
很多中国公司在喀麦隆工作，他们修路、建桥。\\
Hěn duō Zhōngguó gōngsī zài Kāmàilóng gōngzuò, tāmen xiū lù, jiàn qiáo.\\
« Beaucoup de sociétés chinoises travaillent au Cameroun, elles construisent des routes et des ponts. »\\
喀麦隆人也在中国市场买东西，两国的生意越来越多。\\
Kāmàilóng rén yě zài Zhōngguó shìchǎng mǎi dōngxi, liǎng guó de shēngyi yuè lái yuè duō.\\
« Les Camerounais achètent aussi des produits sur les marchés chinois, et le commerce entre les deux pays augmente de plus en plus. »
\end{textebox}

\begin{exoresolu}[Questions de compréhension sur un texte économique]
Énoncé. Lis le texte « 中国经济和喀麦隆 » ci-dessus, puis réponds en chinois aux questions : (1) 中国公司在哪里工作？ (2) 他们做什么？ (3) 为什么两国的生意越来越多？
\tcblower
\textbf{Solution détaillée.}
\begin{enumerate}
\item La question porte sur le lieu (\zh{哪里}) : le texte dit \zh{在喀麦隆工作}. Réponse : \zh{中国公司在喀麦隆工作。} (Zhōngguó gōngsī zài Kāmàilóng gōngzuò.) « Les sociétés chinoises travaillent au Cameroun. » On garde \zh{在} + lieu \textbf{avant} le verbe.
\item \zh{做什么} demande l'action : \zh{他们修路、建桥。} (Tāmen xiū lù, jiàn qiáo.) « Ils construisent des routes et des ponts. » La virgule de liste chinoise \zh{、} sépare les deux actions.
\item \zh{为什么} exige \zh{因为} : \zh{因为很多中国公司在喀麦隆工作，喀麦隆人也在中国市场买东西。} (Yīnwèi hěn duō Zhōngguó gōngsī zài Kāmàilóng gōngzuò, Kāmàilóng rén yě zài Zhōngguó shìchǎng mǎi dōngxi.) « Parce que beaucoup de sociétés chinoises travaillent au Cameroun et que les Camerounais achètent aussi sur les marchés chinois. »
\end{enumerate}
\textbf{Conclusion.} Chaque réponse reprend la structure de la question et cite fidèlement le texte : c'est la méthode la plus sûre pour ne pas perdre de points.
\end{exoresolu}

\courssub{Savoir-faire 5 : Traduire et produire des phrases sur l'économie (ancrage Cameroun--Chine)}

\begin{methode}[Traduire du français vers le chinois sans calque]
\begin{enumerate}
\item \textbf{Identifie le sujet et le verbe} de la phrase française, puis place-les dans l'ordre chinois : Sujet + Temps/Lieu + Verbe + Complément.
\item \textbf{Traduis les compléments de lieu avec \zh{在}} placé avant le verbe : « il travaille à Douala » → \zh{他在杜阿拉工作。} (et non *\zh{他工作在杜阿拉}).
\item \textbf{Utilise \zh{得} pour le complément de manière} : « les affaires vont bien » → \zh{生意做得很好。} (Shēngyi zuò de hěn hǎo.)
\item \textbf{Choisis le bon classificateur} et le bon adverbe de degré : \zh{很} hěn, \zh{太……了} tài... le (trop), \zh{越来越} yuè lái yuè (de plus en plus).
\end{enumerate}
\end{methode}

\begin{exoresolu}[Traduction guidée : le commerce à Yaoundé]
Énoncé. Traduis en chinois : (1) « Mon père travaille dans une société à Yaoundé. » (2) « Les légumes de ce marché sont de moins en moins chers. » (3) « Elle vend très bien ses produits. »
\tcblower
\textbf{Solution détaillée.}
\begin{enumerate}
\item Sujet \zh{我爸爸} + lieu \zh{在雅温得的一家公司} + verbe \zh{工作}. Le classificateur de \zh{公司} est \zh{家}. Réponse : \zh{我爸爸在雅温得的一家公司工作。} (Wǒ bàba zài Yǎwēndé de yì jiā gōngsī gōngzuò.)
\item « De moins en moins cher » = \zh{越来越便宜} ; le sujet est \zh{这个市场的蔬菜} avec le possessif \zh{的}. Réponse : \zh{这个市场的蔬菜越来越便宜。} (Zhège shìchǎng de shūcài yuè lái yuè piányi.)
\item Complément de manière avec \zh{得} : quand le verbe a un objet, on met l'objet en tête de phrase (thématisation) et on reprend le verbe : \zh{产品} + \zh{卖} + \zh{得} + \zh{很好}. Réponse : \zh{她的产品卖得很好。} (Tā de chǎnpǐn mài de hěn hǎo.) « Ses produits se vendent très bien. »
\end{enumerate}
\textbf{Conclusion.} Trois pièges évités : la place de \zh{在} + lieu avant le verbe, la structure \zh{越来越} + adjectif, et le complément de manière avec \zh{得} (et non \zh{的}).
\end{exoresolu}

\begin{savaistu}[Le coquillage, première monnaie de la Chine]
Dans la Chine ancienne, les coquillages servaient de monnaie d'échange. C'est pourquoi la clé \zh{贝} (bèi, coquillage) apparaît dans presque tous les caractères liés à l'argent : \zh{贵} (guì, cher), \zh{货} (huò, marchandise), \zh{买} (mǎi, acheter), \zh{卖} (mài, vendre), \zh{财} (cái, richesse). Au Cameroun comme en Chine, les marchés sont au cœur de la vie économique quotidienne : le marché Mokolo à Yaoundé ou le marché Mboppi à Douala jouent un rôle comparable à celui des grands marchés chinois, où l'on négocie les prix — en chinois, marchander se dit \zh{讲价} (jiǎng jià).
\end{savaistu}

\begin{attention}[Les 5 erreurs qui coûtent des points en classe et à l'examen]
\begin{enumerate}
\item \textbf{Confondre \zh{买} (mǎi, acheter) et \zh{卖} (mài, vendre).} Sons presque identiques, tons différents (3\textsuperscript{e} contre 4\textsuperscript{e}) et caractères différents (\zh{卖} a un \zh{十} en haut). Une inversion change complètement le sens de la phrase.
\item \textbf{Placer le complément de lieu après le verbe} (calque du français) : *\zh{他工作在杜阿拉} est faux. Correct : \zh{他在杜阿拉工作。} La structure est Sujet + \zh{在} + lieu + verbe.
\item \textbf{Confondre \zh{的}, \zh{得} et \zh{地}.} \zh{得} sert au complément de manière après le verbe : \zh{生意做得很好} ; \zh{的} marque la possession ou la qualification : \zh{很便宜的东西}. Écrire *\zh{生意做的很好} est une faute classique.
\item \textbf{Oublier ou mal choisir le classificateur} : on dit \zh{一家公司} (une société), \zh{一家商店} (un magasin), \zh{一个市场} (un marché). Ne jamais écrire *\zh{一公司}.
\item \textbf{Négliger les tons dans le pinyin} : \zh{jīngjì} (\zh{经济}, économie) porte un 1\textsuperscript{er} ton puis un 4\textsuperscript{e} ton ; \zh{shìchǎng} (\zh{市场}, marché) un 4\textsuperscript{e} puis un 3\textsuperscript{e} ton. Un ton faux ou absent est compté comme une erreur, même si le caractère est juste.
\end{enumerate}
\end{attention}

\newpage

\coursec{Exercices}

\courssub{Je vérifie mes connaissances}

\begin{exercice}[QCM : le bon mot]
Pour chaque phrase, choisis la bonne réponse (A, B ou C). Recopie la phrase complète avec le mot choisi.
\begin{enumerate}
\item 喀麦隆的\zh{经济}发展得很快。Kāmàilóng de jīngjì fāzhǎn de hěn kuài.\\
A. 经济 (jīngjì, économie) \quad B. 市场 (shìchǎng, marché) \quad C. 钱 (qián, argent)
\item 我妈妈每天都在\zh{市场}买东西。Wǒ māma měitiān dōu zài \_\_\_ mǎi dōngxi.\\
A. 公司 (gōngsī, entreprise) \quad B. 市场 (shìchǎng, marché) \quad C. 银行 (yínháng, banque)
\item 他做生意，赚了很多\zh{钱}。Tā zuò shēngyi, zhuàn le hěn duō \_\_\_.\\
A. 钱 (qián, argent) \quad B. 工作 (gōngzuò, travail) \quad C. 贵 (guì, cher)
\item 中国是一个很大的\zh{国家}。Zhōngguó shì yí ge hěn dà de \_\_\_.\\
A. 商店 (shāngdiàn, magasin) \quad B. 国家 (guójiā, pays) \quad C. 工厂 (gōngchǎng, usine)
\end{enumerate}
\end{exercice}

\begin{exercice}[Vrai ou faux ?]
Lis chaque phrase et dis si elle est \textbf{vraie} (对 duì) ou \textbf{fausse} (不对 bú duì). Justifie ta réponse en français en une phrase.
\begin{enumerate}
\item \zh{经济} (jīngjì) signifie « économie ».
\item \zh{便宜} (piányi) signifie « cher ».
\item \zh{卖} (mài) est un verbe qui signifie « vendre ».
\item \zh{生意} (shēngyi) signifie « commerce, affaires ».
\item \zh{越来越贵} (yuè lái yuè guì) signifie « de moins en moins cher ».
\end{enumerate}
\end{exercice}

\begin{exercice}[Vocabulaire : complète le tableau]
Recopie et complète le tableau suivant : pour chaque caractère, écris le pinyin (avec les tons) et la traduction française.
\begin{center}
\begin{tabularx}{\linewidth}{|X|X|X|}\hline
\rowcolor{popblueL} Caractères & Pinyin & Traduction \\ \hline
经济 & \ldots\ldots & \ldots\ldots \\ \hline
发展 & \ldots\ldots & \ldots\ldots \\ \hline
市场 & \ldots\ldots & \ldots\ldots \\ \hline
公司 & \ldots\ldots & \ldots\ldots \\ \hline
工作 & \ldots\ldots & \ldots\ldots \\ \hline
买 & \ldots\ldots & \ldots\ldots \\ \hline
\end{tabularx}
\end{center}
\end{exercice}

\begin{exercice}[Associe le pinyin aux caractères]
Relie chaque caractère ou mot de la colonne A à son pinyin de la colonne B.
\begin{center}
\begin{tabularx}{\linewidth}{|X|X|}\hline
\rowcolor{popblueL} Colonne A (caractères) & Colonne B (pinyin) \\ \hline
1. 钱 & a. shāngdiàn \\ \hline
2. 商店 & b. qián \\ \hline
3. 贵 & c. gōngrén \\ \hline
4. 工人 & d. guì \\ \hline
5. 做生意 & e. zuò shēngyi \\ \hline
\end{tabularx}
\end{center}
\end{exercice}

\courssub{Je m'entraîne}

\begin{exercice}[Traduis en français]
Traduis les phrases suivantes en français. Attention à la place des compléments de lieu et de temps.
\begin{enumerate}
\item \zh{喀麦隆的经济发展得很快。} Kāmàilóng de jīngjì fāzhǎn de hěn kuài.
\item \zh{很多人在市场买东西。} Hěn duō rén zài shìchǎng mǎi dōngxi.
\item \zh{他的公司在杜阿拉。} Tā de gōngsī zài Dù'ālā.
\item \zh{中国工人在雅温得工作。} Zhōngguó gōngrén zài Yǎwēndé gōngzuò.
\end{enumerate}
\end{exercice}

\begin{exercice}[Traduis en chinois]
Traduis les phrases suivantes en chinois (caractères obligatoires, pinyin en dessous).
\begin{enumerate}
\item Beaucoup de gens achètent des choses au marché.
\item L'économie se développe de plus en plus vite.
\item Mon père travaille dans une entreprise chinoise.
\item Il fait du commerce et gagne de l'argent.
\end{enumerate}
\end{exercice}

\begin{exercice}[Transforme avec 越来越]
Transforme chaque phrase avec la structure \zh{越来越} + adjectif, comme dans le modèle.\\
\textbf{Modèle :} 生活好。→ \zh{生活越来越好。} Shēnghuó yuè lái yuè hǎo. (La vie devient de plus en plus bonne.)
\begin{enumerate}
\item 公司多。
\item 东西贵。
\item 经济发展快。
\item 学汉语的学生多。
\end{enumerate}
\end{exercice}

\begin{exercice}[Texte à trous]
Complète le texte suivant avec les mots de la liste : \zh{经济}、\zh{市场}、\zh{发展}、\zh{生意}、\zh{越来越}。Recopie le texte complet.
\begin{textebox}[Texte à trous]
喀麦隆的\_\_\_\_发展得很快。Kāmàilóng de \_\_\_\_ fāzhǎn de hěn kuài.\\
在杜阿拉和雅温得，\_\_\_\_上的人很多。Zài Dù'ālā hé Yǎwēndé, \_\_\_\_ shang de rén hěn duō.\\
很多人做\_\_\_\_，卖水果、衣服和鱼。Hěn duō rén zuò \_\_\_\_, mài shuǐguǒ, yīfu hé yú.\\
现在喀麦隆和中国的贸易\_\_\_\_多。Xiànzài Kāmàilóng hé Zhōngguó de màoyì \_\_\_\_ duō.\\
两国的经济都在\_\_\_\_。Liǎng guó de jīngjì dōu zài \_\_\_\_.
\end{textebox}
\end{exercice}

\begin{exercice}[Remets les mots dans l'ordre]
Remets les mots dans le bon ordre pour faire des phrases correctes. Écris chaque phrase en caractères avec le pinyin.
\begin{enumerate}
\item 经济 / 喀麦隆 / 的 / 很快 / 发展得
\item 在 / 东西 / 买 / 市场 / 人们
\item 越来越 / 汉语 / 难 / 吗 / ？
\item 他 / 生意 / 做 / 在 / 杜阿拉
\end{enumerate}
\end{exercice}

\courssub{Je me prépare au Bac}

\begin{exercice}[Compréhension écrite — type Bac]
Lis le texte suivant, puis réponds aux questions \textbf{en chinois, par des phrases complètes}.
\begin{textebox}[Le marché de Mokolo à Yaoundé]
在雅温得有一个很有名的市场，叫莫科洛市场。Zài Yǎwēndé yǒu yí ge hěn yǒumíng de shìchǎng, jiào Mòkēluò shìchǎng.\\
这个市场很大，每天有几千人。Zhège shìchǎng hěn dà, měitiān yǒu jǐ qiān rén.\\
人们在那里买东西：水果、蔬菜、鱼、衣服……Rénmen zài nàli mǎi dōngxi: shuǐguǒ, shūcài, yú, yīfu……\\
很多东西不太贵，所以大家都喜欢去。Hěn duō dōngxi bú tài guì, suǒyǐ dàjiā dōu xǐhuan qù.\\
市场里也有很多中国商品。Shìchǎng li yě yǒu hěn duō Zhōngguó shāngpǐn.\\
现在，喀麦隆和中国的贸易越来越多，Xiànzài, Kāmàilóng hé Zhōngguó de màoyì yuè lái yuè duō,\\
两国的经济发展得越来越快。liǎng guó de jīngjì fāzhǎn de yuè lái yuè kuài.
\end{textebox}
\textbf{Questions :}
\begin{enumerate}
\item 莫科洛市场在哪儿？(Où se trouve le marché de Mokolo ?)
\item 人们在这个市场买什么东西？(Qu'est-ce que les gens achètent dans ce marché ?)
\item 为什么大家都喜欢去这个市场？(Pourquoi tout le monde aime-t-il aller à ce marché ?)
\item 市场里有什么国家的商品？(De quel pays viennent les produits qu'on trouve aussi dans le marché ?)
\item 喀麦隆和中国的贸易怎么样？(Comment évolue le commerce entre le Cameroun et la Chine ?)
\end{enumerate}
\textbf{Question de langue :} relève dans le texte deux phrases avec la structure \zh{越来越} et traduis-les en français.
\end{exercice}

\begin{exercice}[Thème et version — type Bac]
\textbf{A. Thème.} Traduis en chinois (caractères obligatoires) :
\begin{enumerate}
\item L'économie du Cameroun se développe vite.
\item Ma mère achète des légumes au marché tous les jours.
\item Il y a de plus en plus d'entreprises chinoises à Douala.
\item Les produits ne sont pas chers, alors les gens sont contents.
\end{enumerate}
\textbf{B. Version.} Traduis en français :
\begin{textebox}[Version]
中国的经济发展得很快。Zhōngguó de jīngjì fāzhǎn de hěn kuài.\\
很多中国公司来非洲做生意。Hěn duō Zhōngguó gōngsī lái Fēizhōu zuò shēngyi.\\
他们在喀麦隆工作，也帮助喀麦隆发展经济。Tāmen zài Kāmàilóng gōngzuò, yě bāngzhù Kāmàilóng fāzhǎn jīngjì.
\end{textebox}
\end{exercice}

\begin{exercice}[Expression écrite — type Bac]
\textbf{Sujet :} 喀麦隆和中国的贸易 (Le commerce entre le Cameroun et la Chine).\\
Rédige un paragraphe de \textbf{60 à 80 caractères chinois} (environ 5 phrases) sur ce sujet. Tu dois employer \textbf{au moins 8 mots du glossaire} de la leçon (par exemple : 经济、发展、市场、生意、公司、贸易、商品、买、卖、越来越) et \textbf{au moins une fois} la structure \zh{越来越}.\\
\textbf{Consignes :} écris en caractères simplifiés ; soigne l'ordre des mots (sujet + temps/lieu + verbe) ; n'oublie pas la ponctuation chinoise ，。
\end{exercice}

\newpage

\coursec{Corrigés des exercices}

\begin{corrige}[Exercice 1 — QCM : le bon mot]
\textbf{Réponses :}
\begin{enumerate}
\item \textbf{A.} 喀麦隆的\zh{经济}发展得很快。Kāmàilóng de jīngjì fāzhǎn de hěn kuài. « L'économie du Cameroun se développe très vite. » Le sujet du verbe \zh{发展} est naturellement \zh{经济}.
\item \textbf{B.} 我妈妈每天都在\zh{市场}买东西。Wǒ māma měitiān dōu zài shìchǎng mǎi dōngxi. « Ma mère achète des choses au marché tous les jours. » On achète des choses (\zh{买东西}) au marché, pas dans une entreprise ni dans une banque.
\item \textbf{A.} 他做生意，赚了很多\zh{钱}。Tā zuò shēngyi, zhuàn le hěn duō qián. « Il fait du commerce et a gagné beaucoup d'argent. » Le verbe \zh{赚} (gagner) se construit avec \zh{钱}.
\item \textbf{B.} 中国是一个很大的\zh{国家}。Zhōngguó shì yí ge hěn dà de guójiā. « La Chine est un très grand pays. » La Chine est un pays (\zh{国家}).
\end{enumerate}
\textbf{Point de méthode :} dans un QCM de vocabulaire, vérifie toujours la \emph{collocation} : quel mot va naturellement avec le verbe de la phrase (\zh{发展经济}, \zh{买东西}, \zh{赚钱}).
\end{corrige}

\begin{corrige}[Exercice 2 — Vrai ou faux ?]
\begin{enumerate}
\item \textbf{Vrai (对).} \zh{经济} (jīngjì) signifie bien « économie » ; c'est le mot-clé de la leçon.
\item \textbf{Faux (不对).} \zh{便宜} (piányi) signifie « bon marché, pas cher » ; « cher » se dit \zh{贵} (guì). Attention à ce faux ami très fréquent.
\item \textbf{Vrai (对).} \zh{卖} (mài) est le verbe « vendre » ; son contraire est \zh{买} (mǎi, acheter) — seul le ton change : 3\up{e} ton pour acheter, 4\up{e} ton pour vendre.
\item \textbf{Vrai (对).} \zh{生意} (shēngyi) signifie « commerce, affaires » ; \zh{做生意} = « faire du commerce ».
\item \textbf{Faux (不对).} \zh{越来越贵} signifie « de plus en plus cher ». « De moins en moins cher » se dirait \zh{越来越便宜}.
\end{enumerate}
\textbf{Point de vigilance :} la structure \zh{越来越} + adjectif exprime toujours une \emph{augmentation} du degré, jamais une diminution.
\end{corrige}

\begin{corrige}[Exercice 3 — Vocabulaire : complète le tableau]
\begin{center}
\begin{tabularx}{\linewidth}{|X|X|X|}\hline
\rowcolor{popblueL} Caractères & Pinyin & Traduction \\ \hline
经济 & jīngjì & économie (nom) \\ \hline
发展 & fāzhǎn & se développer ; développement \\ \hline
市场 & shìchǎng & marché \\ \hline
公司 & gōngsī & entreprise, société \\ \hline
工作 & gōngzuò & travailler ; travail \\ \hline
买 & mǎi & acheter \\ \hline
\end{tabularx}
\end{center}
\textbf{Points de vigilance :} le ton de \zh{买} est le 3\up{e} ton (mǎi), à ne pas confondre avec \zh{卖} (mài, 4\up{e} ton, vendre). Dans \zh{公司}, les deux caractères sont au 1\up{er} ton : gōngsī.
\end{corrige}

\begin{corrige}[Exercice 4 — Associe le pinyin aux caractères]
\textbf{Correspondances :} 1 — b ; 2 — a ; 3 — d ; 4 — c ; 5 — e.
\begin{itemize}
\item 1. \zh{钱} = qián (argent)
\item 2. \zh{商店} = shāngdiàn (magasin)
\item 3. \zh{贵} = guì (cher)
\item 4. \zh{工人} = gōngrén (ouvrier)
\item 5. \zh{做生意} = zuò shēngyi (faire du commerce)
\end{itemize}
\textbf{Astuce :} repère les syllabes communes : \zh{商店} et \zh{做生意} contiennent tous deux \zh{商} (shāng), le caractère du commerce.
\end{corrige}

\begin{corrige}[Exercice 5 — Traduis en français]
\begin{enumerate}
\item \zh{喀麦隆的经济发展得很快。} « L'économie du Cameroun se développe très vite. » Structure : verbe + \zh{得} + adverbe (complément de manière/degré).
\item \zh{很多人在市场买东西。} « Beaucoup de gens achètent des choses au marché. » En chinois, le complément de lieu \zh{在市场} se place \emph{avant} le verbe ; en français, « au marché » se place après.
\item \zh{他的公司在杜阿拉。} « Son entreprise est à Douala / se trouve à Douala. » Ici \zh{在} est le verbe principal « être situé ».
\item \zh{中国工人在雅温得工作。} « Les ouvriers chinois travaillent à Yaoundé. » Même remarque : le lieu précède le verbe en chinois, il le suit en français.
\end{enumerate}
\textbf{Point de vigilance :} ne traduis jamais mot à mot dans l'ordre chinois ; remets le complément de lieu à sa place française.
\end{corrige}

\begin{corrige}[Exercice 6 — Traduis en chinois]
\begin{enumerate}
\item \zh{很多人在市场买东西。} Hěn duō rén zài shìchǎng mǎi dōngxi. (Le lieu \zh{在市场} avant le verbe \zh{买}.)
\item \zh{经济发展得越来越快。} Jīngjì fāzhǎn de yuè lái yuè kuài. (« De plus en plus vite » = \zh{越来越快} ; ne pas oublier \zh{得} après le verbe.)
\item \zh{我爸爸在一家中国公司工作。} Wǒ bàba zài yì jiā Zhōngguó gōngsī gōngzuò. (Classificateur \zh{家} devant \zh{公司} ; lieu avant le verbe.)
\item \zh{他做生意，赚了很多钱。} Tā zuò shēngyi, zhuàn le hěn duō qián. (\zh{了} marque l'accompli : « a gagné ».)
\end{enumerate}
\textbf{Points de vigilance :} ordre Sujet + lieu + verbe ; \zh{得} pour le complément de degré ; classificateur \zh{家} pour les entreprises.
\end{corrige}

\begin{corrige}[Exercice 7 — Transforme avec 越来越]
\begin{enumerate}
\item \zh{公司越来越多。} Gōngsī yuè lái yuè duō. « Il y a de plus en plus d'entreprises. »
\item \zh{东西越来越贵。} Dōngxi yuè lái yuè guì. « Les choses sont de plus en plus chères. »
\item \zh{经济发展得越来越快。} Jīngjì fāzhǎn de yuè lái yuè kuài. « L'économie se développe de plus en plus vite. » Ici, l'adjectif qualifie le verbe \zh{发展} : on garde donc la structure verbe + \zh{得} + \zh{越来越} + adjectif.
\item \zh{学汉语的学生越来越多。} Xué Hànyǔ de xuésheng yuè lái yuè duō. « Il y a de plus en plus d'élèves qui apprennent le chinois. »
\end{enumerate}
\textbf{Règle :} si l'adjectif est le prédicat direct (phrase 1, 2, 4) : Sujet + \zh{越来越} + adjectif. Si l'adjectif qualifie un verbe d'action (phrase 3) : Sujet + verbe + \zh{得} + \zh{越来越} + adjectif.
\end{corrige}

\begin{corrige}[Exercice 8 — Texte à trous]
\textbf{Texte complété :}
\begin{itemize}
\item 喀麦隆的\zh{经济}发展得很快。Kāmàilóng de jīngjì fāzhǎn de hěn kuài.
\item 在杜阿拉和雅温得，\zh{市场}上的人很多。Zài Dù'ālā hé Yǎwēndé, shìchǎng shang de rén hěn duō.
\item 很多人做\zh{生意}，卖水果、衣服和鱼。Hěn duō rén zuò shēngyi, mài shuǐguǒ, yīfu hé yú.
\item 现在喀麦隆和中国的贸易\zh{越来越}多。Xiànzài Kāmàilóng hé Zhōngguó de màoyì yuè lái yuè duō.
\item 两国的经济都在\zh{发展}。Liǎng guó de jīngjì dōu zài fāzhǎn.
\end{itemize}
\textbf{Traduction :} « L'économie du Cameroun se développe très vite. À Douala et à Yaoundé, il y a beaucoup de monde sur les marchés. Beaucoup de gens font du commerce : ils vendent des fruits, des vêtements et du poisson. Aujourd'hui, le commerce entre le Cameroun et la Chine est de plus en plus important. Les économies des deux pays sont en train de se développer. »\\
\textbf{Indices de résolution :} \zh{做} appelle \zh{生意} ; \zh{越来越} + adjectif (\zh{多}) ; \zh{在} + verbe appelle un verbe d'action, donc \zh{发展}.
\end{corrige}

\begin{corrige}[Exercice 9 — Remets les mots dans l'ordre]
\begin{enumerate}
\item \zh{喀麦隆的经济发展得很快。} Kāmàilóng de jīngjì fāzhǎn de hěn kuài. (Sujet \zh{喀麦隆的经济} + verbe \zh{发展} + \zh{得} + degré.)
\item \zh{人们在市场买东西。} Rénmen zài shìchǎng mǎi dōngxi. (Sujet + lieu \zh{在市场} + verbe.)
\item \zh{汉语越来越难吗？} Hànyǔ yuè lái yuè nán ma? « Le chinois devient-il de plus en plus difficile ? » (Question avec \zh{吗} à la fin.)
\item \zh{他在杜阿拉做生意。} Tā zài Dù'ālā zuò shēngyi. « Il fait du commerce à Douala. » (Le lieu avant le verbe \zh{做}.)
\end{enumerate}
\textbf{Point de vigilance :} l'erreur typique du francophone est de placer le lieu après le verbe (\zh{买东西在市场}) : c'est fautif en chinois.
\end{corrige}

\begin{corrige}[Exercice 10 — Compréhension écrite — type Bac]
\textbf{Réponses en phrases complètes :}
\begin{enumerate}
\item \zh{莫科洛市场在雅温得。} Mòkēluò shìchǎng zài Yǎwēndé. « Le marché de Mokolo se trouve à Yaoundé. »
\item \zh{人们在这个市场买水果、蔬菜、鱼和衣服。} Rénmen zài zhège shìchǎng mǎi shuǐguǒ, shūcài, yú hé yīfu. « Les gens y achètent des fruits, des légumes, du poisson et des vêtements. »
\item \zh{因为很多东西不太贵，所以大家都喜欢去。} Yīnwèi hěn duō dōngxi bú tài guì, suǒyǐ dàjiā dōu xǐhuan qù. « Parce que beaucoup de choses ne sont pas trop chères, tout le monde aime y aller. » (Penser à la paire \zh{因为……所以……}.)
\item \zh{市场里也有很多中国商品。} Shìchǎng li yě yǒu hěn duō Zhōngguó shāngpǐn. « Il y a aussi beaucoup de produits chinois dans le marché. »
\item \zh{喀麦隆和中国的贸易越来越多。} Kāmàilóng hé Zhōngguó de màoyì yuè lái yuè duō. « Le commerce entre le Cameroun et la Chine est de plus en plus important. »
\end{enumerate}
\textbf{Question de langue :} les deux phrases avec \zh{越来越} sont :
\begin{itemize}
\item \zh{喀麦隆和中国的贸易越来越多} : « le commerce entre le Cameroun et la Chine est de plus en plus important » ;
\item \zh{两国的经济发展得越来越快} : « les économies des deux pays se développent de plus en plus vite ».
\end{itemize}
\textbf{Barème indicatif (10 pts) :} 1,5 pt par question de compréhension (7,5 pts) ; 2,5 pts pour la question de langue (relevé 1 pt, traductions 1,5 pt). Toute réponse doit être une phrase complète avec sujet et verbe ; une réponse d'un seul mot perd la moitié des points.
\end{corrige}

\begin{corrige}[Exercice 11 — Thème et version — type Bac]
\textbf{A. Thème.}
\begin{enumerate}
\item \zh{喀麦隆的经济发展得很快。} Kāmàilóng de jīngjì fāzhǎn de hěn kuài.
\item \zh{我妈妈每天在市场买蔬菜。} Wǒ māma měitiān zài shìchǎng mǎi shūcài. (Le temps \zh{每天} et le lieu \zh{在市场} avant le verbe.)
\item \zh{杜阿拉有越来越多的中国公司。} Dù'ālā yǒu yuè lái yuè duō de Zhōngguó gōngsī. (Construction avec \zh{有} ; \zh{越来越多的} + nom.)
\item \zh{商品不贵，所以人们很高兴。} Shāngpǐn bù guì, suǒyǐ rénmen hěn gāoxìng. (Négation \zh{不} devant l'adjectif \zh{贵} ; \zh{所以} introduit la conséquence.)
\end{enumerate}
\textbf{B. Version.} « L'économie chinoise se développe très vite. Beaucoup d'entreprises chinoises viennent faire du commerce en Afrique. Elles travaillent au Cameroun et aident aussi le Cameroun à développer son économie. »\\
\textbf{Barème indicatif (10 pts) :} thème 6 pts (1,5 pt par phrase : 0,5 vocabulaire, 0,5 ordre des mots, 0,5 caractères) ; version 4 pts (fidélité 2 pts, qualité du français 2 pts).\\
\textbf{Points de vigilance :} place de \zh{每天} et du complément de lieu avant le verbe ; \zh{得} dans le complément de degré ; ne pas écrire \zh{越来越} sans adjectif derrière.
\end{corrige}

\begin{corrige}[Exercice 12 — Expression écrite — type Bac]
\textbf{Proposition de rédaction modèle (environ 70 caractères) :}
\begin{textebox}[Modèle de production]
喀麦隆和中国的贸易越来越多。Kāmàilóng hé Zhōngguó de màoyì yuè lái yuè duō.\\
很多中国公司来喀麦隆做生意。Hěn duō Zhōngguó gōngsī lái Kāmàilóng zuò shēngyi.\\
在杜阿拉和雅温得的市场，人们买中国商品。Zài Dù'ālā hé Yǎwēndé de shìchǎng, rénmen mǎi Zhōngguó shāngpǐn.\\
这些商品不太贵，大家都喜欢。Zhèxiē shāngpǐn bú tài guì, dàjiā dōu xǐhuan.\\
两国的经济发展得越来越快。Liǎng guó de jīngjì fāzhǎn de yuè lái yuè kuài.
\end{textebox}
\textbf{Traduction :} « Le commerce entre le Cameroun et la Chine est de plus en plus important. Beaucoup d'entreprises chinoises viennent faire des affaires au Cameroun. Sur les marchés de Douala et de Yaoundé, les gens achètent des produits chinois. Ces produits ne sont pas trop chers et tout le monde les aime. Les économies des deux pays se développent de plus en plus vite. »\\
\textbf{Mots du glossaire employés :} 贸易、公司、做生意、市场、买、商品、经济、发展、越来越 (9 mots + la structure imposée : consigne largement remplie).\\
\textbf{Barème indicatif (10 pts) :} respect du nombre de caractères et des consignes (2 pts) ; emploi d'au moins 8 mots du glossaire (2 pts) ; structure \zh{越来越} correcte (1 pt) ; exactitude grammaticale et ordre des mots (3 pts) ; qualité des caractères et de la ponctuation chinoise (2 pts).\\
\textbf{Points de vigilance :} complément de lieu avant le verbe ; \zh{得} après \zh{发展} ; ponctuation pleine chasse ，。 ; ne pas traduire mot à mot depuis le français.
\end{corrige}

\newpage

\coursec{Fiche bilan}

\coursec{Leçon 31 — Vocabulaire et compréhension de texte : économie}

\courssub{Mise en route et découverte du thème}
\begin{experience}[Observation d'images : le marché et le port]
Observez les deux photos projetées : \textbf{Photo A} : le Marché Central de Yaoundé, étals colorés, négociation des prix (\zh{价格} jiàgé), vendeurs (\zh{商人} shāngrén) et clients. \textbf{Photo B} : le Port autonome de Douala, conteneurs (\zh{集装箱} jízhuāngxiāng), grues, chargement de bois (\zh{木材} mùcái) et coton (\zh{棉花} miánhuā) vers des navires pour la Chine.
\par\medskip
\textbf{Activité orale (paires) :} En vous aidant du lexique ci-dessous, décrivez chaque image en 2 phrases chinoises simples.
\begin{center}
\begin{tabularx}{\linewidth}{|X|X|X|X|}
\hline
\rowcolor{popblueL} Caractères & Pinyin & Nature & Traduction \\
\hline
市场 & shìchǎng & n. & marché \\
\hline
港口 & gǎngkǒu & n. & port \\
\hline
出口 & chūkǒu & v. & exporter \\
\hline
进口 & jìnkǒu & v. & importer \\
\hline
产品 & chǎnpǐn & n. & produit \\
\hline
贵 / 便宜 & guì / piányi & adj. & cher / bon marché \\
\hline
\end{tabularx}
\end{center}
\textbf{Production attendue :} \zh{这是市场。价格不贵。} (C'est le marché. Les prix ne sont pas chers.) / \zh{那是港口。很多产品出口。} (C'est le port. Beaucoup de produits sont exportés.)
\end{experience}

\courssub{Texte support : lecture et compréhension globale}
\begin{textebox}[Texte original : La coopération économique Cameroun--Chine]
\zh{中喀经济贸易合作} \\
\zh{Zhōng-Kā jīngjì màoyì hézuò} \\
\zh{喀麦隆和中国的经济合作越来越重要。} \\
\zh{Kāmàilóng hé Zhōngguó de jīngjì hézuò yuèláiyuè zhòngyào.} \\
\zh{很多喀麦隆的农产品，比如木材和棉花，从杜阿拉港出口到中国。} \\
\zh{Hěn duō Kāmàilóng de nóngchǎnpǐn, bǐrú mùcái hé miánhuā, cóng Dù'ālā gǎng chūkǒu dào Zhōngguó.} \\
\zh{中国的工业产品，比如手机和机器，进口到喀麦隆的市场。} \\
\zh{Zhōngguó de gōngyè chǎnpǐn, bǐrú shǒujī hé jīqì, jìnkǒu dào Kāmàilóng de shìchǎng.} \\
\zh{价格比较便宜，质量也好。} \\
\zh{Jiàgé bǐjiào piányi, zhìliàng yě hǎo.} \\
\zh{在雅温得和杜阿拉，有很多中国公司做生意。} \\
\zh{Zài Yǎwēndé hé Dù'ālā, yǒu hěn duō Zhōngguó gōngsī zuò shēngyi.} \\
\zh{喀麦隆的工人学习中国的技术，经济发展得很快。} \\
\zh{Kāmàilóng de gōngrén xuéxí Zhōngguó de jìshù, jīngjì fāzhǎn de hěn kuài.} \\
\zh{未来，两国的贸易会更多。} \\
\zh{Wèilái, liǎngguó de màoyì huì gèng duō.}

\par\medskip
\textbf{Traduction :} La coopération économique et commerciale sino-camerounaise devient de plus en plus importante. Beaucoup de produits agricoles camerounais, comme le bois et le coton, sont exportés depuis le port de Douala vers la Chine. Les produits industriels chinois, comme les téléphones portables et les machines, sont importés sur le marché camerounais. Les prix sont relativement bas et la qualité est bonne. À Yaoundé et Douala, il y a beaucoup d'entreprises chinoises qui font des affaires. Les ouvriers camerounais apprennent la technologie chinoise, l'économie se développe très vite. À l'avenir, le commerce entre les deux pays sera encore plus important.
\end{textebox}

\courssub{Glossaire du thème : économie}
\begin{definition}[Vocabulaire clé du texte (HSK 2--3)]
\begin{center}
\begin{tabularx}{\linewidth}{|X|X|X|X|}
\hline
\rowcolor{popblueL} Caractères & Pinyin & Nature & Traduction \\
\hline
经济 & jīngjì & n. & économie \\
\hline
发展 & fāzhǎn & v. & développer, développement \\
\hline
合作 & hézuò & v./n. & coopérer, coopération \\
\hline
贸易 & màoyì & n. & commerce, échanges commerciaux \\
\hline
市场 & shìchǎng & n. & marché \\
\hline
价格 & jiàgé & n. & prix \\
\hline
贵 & guì & adj. & cher \\
\hline
便宜 & piányi & adj. & bon marché, pas cher \\
\hline
出口 & chūkǒu & v. & exporter \\
\hline
进口 & jìnkǒu & v. & importer \\
\hline
公司 & gōngsī & n. & société, entreprise \\
\hline
生意 & shēngyi & n. & affaires, commerce \\
\hline
产品 & chǎnpǐn & n. & produit \\
\hline
农产品 & nóngchǎnpǐn & n. & produit agricole \\
\hline
工业产品 & gōngyè chǎnpǐn & n. & produit industriel \\
\hline
木材 & mùcái & n. & bois, bois d'œuvre \\
\hline
棉花 & miánhuā & n. & coton \\
\hline
手机 & shǒujī & n. & téléphone portable \\
\hline
机器 & jīqì & n. & machine \\
\hline
质量 & zhìliàng & n. & qualité \\
\hline
技术 & jìshù & n. & technologie, technique \\
\hline
工人 & gōngrén & n. & ouvrier, travailleur \\
\hline
农民 & nóngmín & n. & paysan \\
\hline
港口 & gǎngkǒu & n. & port \\
\hline
未来 & wèilái & n. & avenir, futur \\
\hline
两国 & liǎngguó & n. & les deux pays \\
\hline
比较 & bǐjiào & adv. & relativement, assez \\
\hline
比如 & bǐrú & conj. & par exemple \\
\hline
从 \dots 到 \dots & cóng \dots dào \dots & structure & de \dots à \dots (origine \rightarrow destination) \\
\hline
\end{tabularx}
\end{center}
\end{definition}

\courssub{Clés (radicaux) et ordre des traits}
\begin{propriete}[Analyse des caractères : radicaux « argent » et « métal »]
\textbf{1. Clé \zh{贝} bèi (coquillage \rightarrow ancienne monnaie)} \\
Cette clé indique un lien avec l'argent, la valeur, les transactions. Elle se place généralement à gauche (\zh{贝}) ou en bas (\zh{贱}).
\begin{center}
\begin{tabularx}{\linewidth}{|X|X|X|X|}
\hline
\rowcolor{poporangeL} Caractère & Pinyin & Sens & Décomposition (ordre traits) \\
\hline
\zh{贝} & bèi & coquillage / radical & 4 traits : horizontale, verticale, horizontale pliée, point \\
\hline
\zh{贵} & guì & cher & \zh{贝} (bas) + \zh{中} (haut) + traits supérieurs. Ordre : haut \rightarrow bas. \\
\hline
\zh{贸} & mào & commerce & \zh{贝} (gauche) + \zh{冒} (droite). Ordre : gauche \rightarrow droite. \\
\hline
\zh{货} & huò & marchandise & \zh{贝} (gauche) + \zh{化} (droite). Ordre : gauche \rightarrow droite. \\
\hline
\zh{买} & mǎi & acheter & \zh{卖} (mǎi) sans le trait du haut ? Non, \zh{买} = \zh{卖} sans le \zh{十} ? Rappel : \zh{卖} mài (vendre) a un trait en plus (le \zh{十} « dix » = prix). \zh{买} mǎi (acheter) n'a pas le « dix ». \\
\hline
\end{tabularx}
\end{center}

\textbf{2. Clé \zh{钅} jīn (métal / or \rightarrow monnaie moderne)} \\
Forme compressée de \zh{金} (or, métal) placée à gauche.
\begin{center}
\begin{tabularx}{\linewidth}{|X|X|X|X|}
\hline
\rowcolor{poporangeL} Caractère & Pinyin & Sens & Décomposition \\
\hline
\zh{钅} / \zh{金} & jīn & métal / or & 5 traits (forme gauche) / 8 traits (forme pleine) \\
\hline
\zh{钱} & qián & argent, monnaie & \zh{钅} + \zh{戋} (clé du métal + élément phonétique) \\
\hline
\end{tabularx}
\end{center}
\end{propriete}


\findecours

\end{document}
